\documentclass[11pt]{article}
\usepackage[T1]{fontenc}
\usepackage[margin=1in]{geometry}
\usepackage{amsmath,amssymb,amsthm,microtype}
\usepackage{mathrsfs}
\usepackage{tikz}
\usetikzlibrary{shapes.geometric,arrows.meta,calc}
\usepackage{caption}
\captionsetup{font=small,labelfont=bf,hypcap=false}
\usepackage{placeins}
\renewcommand{\topfraction}{0.85}
\renewcommand{\textfraction}{0.1}
\renewcommand{\floatpagefraction}{0.8}
\widowpenalty=10000
\clubpenalty=10000
\usepackage[hidelinks]{hyperref}
\hypersetup{pdftitle={A Geometric View of Einstein-Yang-Mills Energy and Stress},pdfauthor={Sina Montazeri}}
\newtheorem{theorem}{Theorem}
\newtheorem{proposition}{Proposition}
\newtheorem{lemma}{Lemma}
\theoremstyle{definition}
\newtheorem{definition}{Definition}
% Canonical equations and notation. Imported by main.tex.
% Active semantic definitions for the source-derived manuscript.
% Explanatory underbraces and uniform boxes for principal equations.
\newcommand{\ExplainTerm}[2]{\underbrace{#1}_{\text{\scriptsize #2}}}
\newcommand{\KeyMath}[1]{\boxed{\begin{gathered}#1\end{gathered}}}


\newcommand{\PhaseGroup}{\mathrm{U}(1)}
\newcommand{\PairGroup}{\mathrm{SU}(2)}
\newcommand{\ColorGroup}{\mathrm{SU}(3)}
\newcommand{\ColorAlgebra}{\mathfrak{su}(3)}
\newcommand{\SourceCoupling}{\Lambda}


\newcommand{\SourcePlaneBasis}{(e_r,e_\perp)}

\newcommand{\ComparisonSpace}{{\mathcal V}}
\newcommand{\CouplingMap}{\Lambda}
\newcommand{\SourceFamily}{X}
\newcommand{\SourcePairing}{\langle\,\cdot\,,\,\cdot\,\rangle_{\mathrm F}}
\newcommand{\PlaneInjection}{{\mathsf V}\colon\mathbb R^2\to {\mathcal V}}
\newcommand{\TwoColumnFrame}{{\mathsf V}}


\newcommand{\FirstColumn}{{\mathsf V_0}}
\newcommand{\NormalColumn}{{\mathsf V_1}}


\newcommand{\RealFrame}{\mathsf E}
\newcommand{\BlockComplexStructure}{J}


\newcommand{\LastPairGroup}{\mathrm{SU}(2)_{23}}
\newcommand{\MatrixUnit}{E_{ij}}

\newcommand{\DisplaySize}{\ell}

\newcommand{\TotalDefect}{4\pi}

\newcommand{\FamilyIndex}{p}
\newcommand{\PerturbationSymbol}{\delta X}
\newcommand{\VertexPairOne}{\pm e_1}
\newcommand{\VertexLabel}[2]{#2e_{#1}}
\newcommand{\PhaseRotation}{\mathsf R(\theta)}

\newcommand{\UnitaryConnectionEquation}{\nabla=d+i\,a,\qquad F_\nabla=i\,da}
\newcommand{\PhaseGaugeEquation}{a\longmapsto a+d\chi,\qquad da\longmapsto da}
\newcommand{\LoopPhaseEquation}{\exp\!\left(i\oint_{\partial \mathcal S}a\right)=\exp\!\left(i\int_{\mathcal S} da\right)}

\newcommand{\TripleLiftEquation}{\begin{aligned}\iota{}v&=(v,{L_X}v,{L_Y}{L_X}v),\\{\iota_1^\perp}v&=(-{L_X}^*v,v,0),\qquad{\iota_2^\perp}v=(0,-{L_Y}^*v,v),\end{aligned}}
\newcommand{\OrthogonalColumnsEquation}{\begin{gathered}{\mathsf V_0}=\iota{}{\mathsf V},\qquad {\mathsf V_1}={\iota_1^\perp}{\mathsf V},\qquad {\mathsf W_2}={\iota_2^\perp}{\mathsf V},\\{\mathsf V_2}={\mathsf W_2}-{\mathsf V_1}({\mathsf V_1}^*{\mathsf V_1})^{-1}{\mathsf V_1}^*{\mathsf W_2}\end{gathered}}

\newcommand{\CompatibleConnectionEquation}{\Omega_{\mathbb R}=\mathsf E^*d\mathsf E,\qquad A_J=\tfrac12(\Omega_{\mathbb R}-J\Omega_{\mathbb R} J),\qquad A=A_J-\tfrac13\operatorname{tr}_{\mathbb C}(A_J)I_3}
\newcommand{\PairGroupEquation}{\begin{pmatrix}\alpha&\beta\\-\bar\beta&\bar\alpha\end{pmatrix},\qquad |\alpha|^2+|\beta|^2=1}

\newcommand{\GeneratorDefinitionEquation}{\mathsf A_{ij}=E_{ij}-E_{ji},\quad \mathsf B_{ij}=i(E_{ij}+E_{ji}),\quad \mathsf H_{ij}=i(E_{ii}-E_{jj})\quad(i<j)}

\newcommand{\CrossPairEquation}{[\mathsf A_{12},\mathsf A_{23}]=\mathsf A_{13},\qquad[\mathsf A_{12},\mathsf B_{23}]=\mathsf B_{13}}
\newcommand{\DiagonalDependenceEquation}{\mathsf H_{13}=\mathsf H_{12}+\mathsf H_{23}}

\newcommand{\DiagonalModesEquation}{\frac{\lambda_3}{\sqrt2}=\frac{E_{11}-E_{22}}{\sqrt2},\qquad \frac{\lambda_8}{\sqrt2}=\frac{E_{11}+E_{22}-2E_{33}}{\sqrt6},\qquad \tfrac12\operatorname{tr}(\lambda_a\lambda_b)=\delta_{ab}\quad(a,b\in\{3,8\})}

\newcommand{\PolynomialCouplingEquation}{\Lambda(X)_p=w_p(X)X_p,\qquad w_p(X)=\sum_{j\ne p}\langle X_p,X_j\rangle_{\mathrm F}}
\newcommand{\CouplingDifferentialEquation}{(D\Lambda(X)\delta X)_p=w_p(X)(\delta X)_p+X_p\sum_{j\ne p}\bigl[\langle (\delta X)_p,X_j\rangle_{\mathrm F}+\langle X_p,(\delta X)_j\rangle_{\mathrm F}\bigr]}
\newcommand{\FullChainRuleEquation}{DY[\eta]={L_X}\,DX[\eta],\quad DZ[\eta]={L_Y}{L_X}\,DX[\eta],\quad d{L_Y}[\eta]=D^2\Lambda(Y)[{L_X}\,DX[\eta],\,\cdot\,]}

%


%

% Explicit field dependence, induced phase connection and compatibility.

\newcommand{\CompatibleFieldCurve}{X(s)}

\newcommand{\CouplingWeightDifferential}{Dw_p(X)[\delta X]}
\newcommand{\SecondPerturbationSymbol}{\delta_2X}
\newcommand{\CouplingSecondDifferentialEquation}{\begin{aligned}
(D^2\Lambda(X)[\delta_1X,\delta_2X])_p={}&Dw_p(X)[\delta_1X](\delta_2X)_p+Dw_p(X)[\delta_2X](\delta_1X)_p\\
&+X_p\sum_{j\ne p}\bigl[\langle(\delta_1X)_p,(\delta_2X)_j\rangle_{\mathrm F}+\langle(\delta_2X)_p,(\delta_1X)_j\rangle_{\mathrm F}\bigr].
\end{aligned}}
\newcommand{\BothDifferentialVariationsEquation}{\begin{aligned}d{L_X}[\eta]&=D^2\Lambda(X)[DX[\eta],\,\cdot\,],\\d{L_Y}[\eta]&=D^2\Lambda(Y)[{L_X}\,DX[\eta],\,\cdot\,],\\d({L_Y}{L_X}v)&=(d{L_Y}){L_X}v+{L_Y}(d{L_X})v+{L_Y}{L_X}\,dv.\end{aligned}}
\newcommand{\CombinedFieldMap}{\Phi(X)}
\newcommand{\CompatibilityEquality}{f(\Phi(X))=0}
\newcommand{\FieldVelocityAcceleration}{\delta X=\dot X(0)}
\newcommand{\PropagatedCompatibilityEquation}{\begin{aligned}Df\,D\Phi\,\delta X&=0,\\Df\bigl(D\Phi\,\ddot X(0)+D^2\Phi[\delta X,\delta X]\bigr)&+D^2f[D\Phi\,\delta X,D\Phi\,\delta X]=0.\end{aligned}}
\newcommand{\CompatibilityInequality}{\varphi(\Phi(X))\ge0}
\newcommand{\ContactDerivativeEquation}{\left.\frac{d}{ds}\varphi(\Phi(X(s)))\right|_{s=0}=D\varphi\,D\Phi\,\delta X\ge0}

\newcommand{\TripleComparisonSpace}{{\mathcal V}\oplus {\mathcal V}\oplus {\mathcal V}}

% Equations for the author-requested two-page Phase/Coupling insert.
\newcommand{\DraftTransverseEquation}{X_{i\perp p}=X_i-\frac{\langle X_p,X_i\rangle_{\mathrm F}}{\langle X_p,X_p\rangle_{\mathrm F}}X_p}
\newcommand{\DraftNormalizationEquation}{{e_r}=\frac{X_p}{\|X_p\|_{\mathrm F}},\qquad {e_\perp}=\frac{X_{i\perp p}}{\|X_{i\perp p}\|_{\mathrm F}},\qquad\langle {e_r},{e_\perp}\rangle_{\mathrm F}=0}
\newcommand{\DraftComplexStructureEquation}{J{e_r}={e_\perp},\qquad J{e_\perp}=-{e_r},\qquad J^2=-I}
\newcommand{\DraftRotationEquation}{\mathsf R(\theta)=\cos\theta\,I+\sin\theta\,J}
\newcommand{\DraftPhaseConnectionEquation}{a=\langle {e_\perp},d e_r\rangle_{\mathrm F}}
\newcommand{\DraftProjectedDerivativesEquation}{(d e_r)_{\parallel}=a\,{e_\perp},\qquad(d e_\perp)_{\parallel}=-a\,{e_r}}
\newcommand{\DraftCouplingEquation}{Y=\Lambda(X),\qquad\Lambda(X)_p=w_p(X)X_p,\qquad w_p(X)=\sum_{j\ne p}\langle X_p,X_j\rangle_{\mathrm F}}
\newcommand{\DraftCouplingResponseEquation}{({L_X}\delta X)_p=w_p(X)(\delta X)_p+X_p\sum_{j\ne p}\left[\langle (\delta X)_p,X_j\rangle_{\mathrm F}+\langle X_p,(\delta X)_j\rangle_{\mathrm F}\right],\qquad {L_X}=D\Lambda(X)}
\newcommand{\DraftPairLiftEquation}{{\iota_{L_X}}v=(v,{L_X}v),\qquad{\iota_{L_X}^{\perp}}w=(-{L_X}^*w,w)}
\newcommand{\DraftPairOrthogonalityEquation}{\begin{aligned}\langle{\iota_{L_X}}v,{\iota_{L_X}^{\perp}}w\rangle&=-\langle v,{L_X}^*w\rangle+\langle {L_X}v,w\rangle\\&=0.\end{aligned}}
\newcommand{\DraftPairColumnsEquation}{{\mathsf V_0}={\iota_{L_X}}{\mathsf V},\qquad {\mathsf V_1}={\iota_{L_X}^{\perp}}{\mathsf V}}
\newcommand{\DraftPairGramEquation}{{\mathsf V_0}^*{\mathsf V_1}=0,\qquad {\mathsf V_0}^*{\mathsf V_0}={\mathsf V}^*(I+{L_X}^*{L_X}){\mathsf V}>0,\qquad {\mathsf V_1}^*{\mathsf V_1}={\mathsf V}^*(I+{L_X}{L_X}^*){\mathsf V}>0}
\newcommand{\DraftNextCoupling}{Z=\Lambda(Y)}

% Equations for the additive Algebra draft.
\newcommand{\DraftThreeLevelVariation}{\delta X\longmapsto(\delta X,{L_X}\delta X,{L_Y}{L_X}\delta X)}
\newcommand{\DraftJointOrthogonality}{{\mathsf V_0}^*{\mathsf V_1}={\mathsf V_0}^*{\mathsf V_2}={\mathsf V_1}^*{\mathsf V_2}=0}
\newcommand{\DraftConnectionMatrix}{A=\begin{pmatrix}ia_1&A_{12}&A_{13}\\-\overline{A_{12}}&ia_2&A_{23}\\-\overline{A_{13}}&-\overline{A_{23}}&ia_3\end{pmatrix},\qquad a_1+a_2+a_3=0}
\newcommand{\DraftGaugeCurvature}{F_A=dA+A\wedge A}
\newcommand{\DraftCrossCurvature}{(F_A)_{13}=dA_{13}+i(a_1-a_3)\wedge A_{13}+A_{12}\wedge A_{23}}

% Equations for the additive Curvature draft.
\newcommand{\DraftSignedEndpoints}{\pm\ell e_1,\qquad\pm\ell e_2,\qquad\pm\ell e_3,\qquad\ell>0}
\newcommand{\DraftOctahedralHull}{P_\ell=\operatorname{conv}\{\pm\ell e_1,\pm\ell e_2,\pm\ell e_3\}}
\newcommand{\DraftOctahedralFaces}{\tau_s=\operatorname{conv}\{\ell s_1e_1,\ell s_2e_2,\ell s_3e_3\},\qquad s\in\{-1,+1\}^{3}}
\newcommand{\DraftClosedSurface}{\partial\!\left(\sum_s \tau_s\right)=0}
\newcommand{\DraftAngularDefect}{\delta_v=2\pi-\sum_{f\ni v}\theta_{vf}}
\newcommand{\DraftRegularDefect}{\delta_v=2\pi-4\frac{\pi}{3}=\frac{2\pi}{3}}
\newcommand{\DraftTotalDefect}{\sum_v\delta_v=4\pi}
\newcommand{\DraftInvariantDefect}{\sum_v\delta_v=2\pi\cdot6-\pi\cdot8=4\pi}

% Equations for the additive Energy draft.

\newcommand{\DraftTotalMassRelation}{M^2c^4=E^2-c^2|\mathbf P|^2,\qquad E_0=E\big|_{\mathbf P=0}=Mc^2}


% Equations for the additive Spectrum draft.


% Scale notation for the consolidated Curvature section.

% Inline field, geometric, energy and spectral notation.
\newcommand{\IntermediateField}{Y}
\newcommand{\SecondCouplingDifferential}{{L_Y}=D\Lambda(Y)}
\newcommand{\CoupledVariation}{{L_X}\delta X}
\newcommand{\FirstAdjacentGroup}{\mathrm{SU}(2)_{12}}
\newcommand{\ColumnIndex}{j}
\newcommand{\RowIndex}{i}
\newcommand{\DiagonalConnectionForms}{a_1,a_2,a_3}
\newcommand{\OffDiagonalConnectionForms}{A_{12},A_{23},A_{13}}
\newcommand{\HermitianConversionFactor}{-i}
\newcommand{\FirstCouplingLabel}{\SourceCoupling}
\newcommand{\SecondCouplingLabel}{\SourceCoupling}
\newcommand{\FirstDifferential}{{L_X}}
\newcommand{\PairComparisonSpace}{{\mathcal V}\oplus {\mathcal V}}
\newcommand{\FirstDifferentialAdjoint}{{L_X}^*}
\newcommand{\FirstCouplingDifferential}{{L_X}=D\Lambda(X)}
\newcommand{\SurfaceVertex}{v}
\newcommand{\VertexFaceAngle}{\theta_{vf}}
\newcommand{\SurfaceFace}{f}

\newcommand{\RadialFieldComponent}{X_p}
\newcommand{\TransverseFieldComponent}{X_i}
\newcommand{\TransverseDirection}{X_{i\perp p}}
\newcommand{\SpanningSurface}{\mathcal S}
\newcommand{\PhaseCurvatureForm}{da}

% Mathematical figure labels; TikZ geometry remains in figures/.
\newcommand{\DisplayDirection}[1]{\ifcase#1\or x\or y\or z\fi}
\newcommand{\CarrierDirectionLabel}[1]{\DisplayDirection{#1}}
\newcommand{\FieldSpatialCoordinates}{(x,y,z)}
\newcommand{\FieldSpatialPhase}{a_{\mathrm{sp}}}
\newcommand{\FieldSpatialAngle}{\chi(x,y,z)}
\newcommand{\FieldSpatialConnectionDisplay}{%
\begin{equation}
a_{\mathrm{sp}}=a_x\,dx+a_y\,dy+a_z\,dz,\qquad
D_j=\partial_j+ia_j,\quad j\in\{x,y,z\}.
\label{eq:one-field-spatial-connection}
\end{equation}%
}
\newcommand{\FieldSpatialGaugeDisplay}{%
\begin{equation}
\widetilde a_{\mathrm{sp}}=a_{\mathrm{sp}}+d\chi,\qquad
\widetilde a_j=a_j+\partial_j\chi,\quad j\in\{x,y,z\}.
\label{eq:one-field-spatial-gauge}
\end{equation}%
}
\newcommand{\FieldSpatialCurvatureDisplay}{%
\begin{equation}
\begin{aligned}
d a_{\mathrm{sp}}={}&(\partial_xa_y-\partial_ya_x)\,dx\wedge dy\\
&+(\partial_ya_z-\partial_za_y)\,dy\wedge dz\\
&+(\partial_za_x-\partial_xa_z)\,dz\wedge dx.
\end{aligned}
\label{eq:one-field-spatial-curvature}
\end{equation}%
}
\newcommand{\FirstPairFigureLabel}{X\xrightarrow{\Lambda}Y\quad\PairGroup_{12}}
\newcommand{\AdjacentPairsFigureLabel}{X\xrightarrow{\Lambda}Y\xrightarrow{\Lambda}Z\quad\PairGroup_{12},\PairGroup_{23}}
\newcommand{\RedAntiGreen}{r\bar g}
\newcommand{\GreenAntiRed}{g\bar r}
\newcommand{\GreenAntiBlue}{g\bar b}
\newcommand{\BlueAntiGreen}{b\bar g}
\newcommand{\RedAntiBlue}{r\bar b}
\newcommand{\BlueAntiRed}{b\bar r}
\newcommand{\FirstDiagonalModeLabel}{\lambda_3/\sqrt2}
\newcommand{\SecondDiagonalModeLabel}{\lambda_8/\sqrt2}
\newcommand{\RedAntiGreenSigns}{+--}
\newcommand{\GreenAntiRedSigns}{-++}
\newcommand{\GreenAntiBlueSigns}{-+-}
\newcommand{\BlueAntiGreenSigns}{+-+}
\newcommand{\RedAntiBlueSigns}{++-}
\newcommand{\BlueAntiRedSigns}{--+}
\newcommand{\FirstDiagonalSigns}{+++}
\newcommand{\SecondDiagonalSigns}{---}

% Complete displayed equations, labels and punctuation in manuscript order.

% Phase
\newcommand{\PhaseTransverseDisplay}{%
\begin{equation}
\begin{gathered}
\ExplainTerm{X_{i\perp p}}{transverse component}
=\ExplainTerm{X_i}{original component}
-\ExplainTerm{\frac{\langle X_p,X_i\rangle_{\mathrm F}}{\langle X_p,X_p\rangle_{\mathrm F}}X_p}{projection along the radial component}
\end{gathered}
\label{eq:radial-tangent}
\end{equation}%
}
\newcommand{\PhaseNormalizationDisplay}{%
\begin{equation}
\DraftNormalizationEquation.
\end{equation}%
}
\newcommand{\PhaseComplexStructureDisplay}{%
\begin{equation}
\DraftComplexStructureEquation.
\label{eq:phase}
\end{equation}%
}
\newcommand{\PhaseRotationDisplay}{%
\begin{equation}
\DraftRotationEquation,
\end{equation}%
}
\newcommand{\PhasePhaseConnectionDisplay}{%
\begin{equation}
\begin{gathered}
\ExplainTerm{a}{phase connection}
=\ExplainTerm{\langle e_\perp,de_r\rangle_{\mathrm F}}{radial change along the transverse direction}
\end{gathered}
\end{equation}%
}
\newcommand{\PhaseProjectedDerivativesDisplay}{%
\begin{equation}
\DraftProjectedDerivativesEquation.
\end{equation}%
}
\newcommand{\PhaseUnitaryConnectionDisplay}{%
\begin{equation}
\UnitaryConnectionEquation.
\label{eq:unitary}
\end{equation}%
}
\newcommand{\PhaseLoopPhaseDisplay}{%
\begin{equation}
\LoopPhaseEquation,
\label{eq:loop}
\end{equation}%
}
\newcommand{\PhasePhaseGaugeDisplay}{%
\begin{equation}
\PhaseGaugeEquation.
\label{eq:phase-gauge}
\end{equation}%
}

% Coupling
\newcommand{\CouplingCouplingDisplay}{%
\begin{equation}
\KeyMath{
\ExplainTerm{Y=\Lambda(X)}{coupled response},\qquad
\ExplainTerm{\Lambda(X)_p=w_p(X)X_p}{component response},\\
\ExplainTerm{w_p(X)=\sum_{j\ne p}\langle X_p,X_j\rangle_{\mathrm F}}{total pairing with the other components}
}
\label{eq:successive}
\end{equation}%
}
\newcommand{\CouplingCouplingResponseDisplay}{%
\begin{equation}
\KeyMath{
\ExplainTerm{L_X=D\Lambda(X)}{derivative at the field},\\
(L_X\delta X)_p=
\ExplainTerm{w_p(X)(\delta X)_p}{component change}
+\ExplainTerm{X_p\sum_{j\ne p}\bigl[\langle(\delta X)_p,X_j\rangle_{\mathrm F}+\langle X_p,(\delta X)_j\rangle_{\mathrm F}\bigr]}{coupling-weight change}
}
\label{eq:coupling-response}
\end{equation}%
}
\newcommand{\CouplingPairLiftDisplay}{%
\begin{equation}
\begin{gathered}
\ExplainTerm{\iota_{L_X}v=(v,L_Xv)}{field and its response},\qquad
\ExplainTerm{\iota_{L_X}^{\perp}w=(-L_X^*w,w)}{perpendicular compensation}
\end{gathered}
\label{eq:pair}
\end{equation}%
}
\newcommand{\CouplingPairOrthogonalityDisplay}{%
\begin{equation}
\DraftPairOrthogonalityEquation
\end{equation}%
}
\newcommand{\CouplingPairColumnsDisplay}{%
\begin{equation}
\DraftPairColumnsEquation.
\end{equation}%
}
\newcommand{\CouplingPairGramDisplay}{%
\begin{equation}
\DraftPairGramEquation.
\end{equation}%
}
\newcommand{\CouplingPairGroupDisplay}{%
\begin{equation}
U=\PairGroupEquation.
\end{equation}%
}

% Algebra
\newcommand{\AlgebraThreeLevelVariationDisplay}{%
\begin{equation}
\KeyMath{
\delta X\longmapsto
(\ExplainTerm{\delta X}{field change},\ExplainTerm{L_X\delta X}{first response},\ExplainTerm{L_YL_X\delta X}{second response})
}
\end{equation}%
}
\newcommand{\AlgebraTripleLiftDisplay}{%
\begin{equation}
\TripleLiftEquation
\label{eq:lifts}
\end{equation}%
}
\newcommand{\AlgebraOrthogonalColumnsDisplay}{%
\begin{equation}
\OrthogonalColumnsEquation.
\label{eq:response-plane-projection}
\end{equation}%
}
\newcommand{\AlgebraJointOrthogonalityDisplay}{%
\begin{equation}
\DraftJointOrthogonality.
\end{equation}%
}
\newcommand{\AlgebraGeneratorDefinitionDisplay}{%
\begin{equation}
\GeneratorDefinitionEquation.
\label{eq:generators}
\end{equation}%
}
\newcommand{\AlgebraCrossPairDisplay}{%
\begin{equation}
\CrossPairEquation.
\label{eq:cross-brackets}
\end{equation}%
}
\newcommand{\AlgebraDiagonalDependenceDisplay}{%
\begin{equation}
\DiagonalDependenceEquation,
\end{equation}%
}
\newcommand{\AlgebraCompatibleConnectionDisplay}{%
\begin{equation}
\CompatibleConnectionEquation.
\label{eq:compatible}
\end{equation}%
}
\newcommand{\AlgebraConnectionMatrixDisplay}{%
\begin{equation}
\DraftConnectionMatrix,
\label{eq:connection}
\end{equation}%
}
\newcommand{\AlgebraGaugeCurvatureDisplay}{%
\begin{equation}
\DraftGaugeCurvature.
\end{equation}%
}
\newcommand{\AlgebraCrossCurvatureDisplay}{%
\begin{equation}
\DraftCrossCurvature.
\label{eq:curvature}
\end{equation}%
}
\newcommand{\AlgebraDiagonalModesDisplay}{%
\begin{equation}
\DiagonalModesEquation.
\label{eq:diagonals}
\end{equation}%
}

% Curvature
\newcommand{\CurvatureSignedEndpointsDisplay}{%
\begin{equation}
\DraftSignedEndpoints.
\end{equation}%
}
\newcommand{\CurvatureOctahedralHullDisplay}{%
\begin{equation}
\DraftOctahedralHull.
\label{eq:octahedron}
\end{equation}%
}
\newcommand{\CurvatureOctahedralFacesDisplay}{%
\begin{equation}
\DraftOctahedralFaces.
\end{equation}%
}
\newcommand{\CurvatureClosedSurfaceDisplay}{%
\begin{equation}
\KeyMath{
\ExplainTerm{\partial\!\left(\sum_s\tau_s\right)}{boundary of the assembled faces}
=\ExplainTerm{0}{shared edges cancel}
}
\end{equation}%
}
\newcommand{\CurvatureAngularDefectDisplay}{%
\begin{equation}
\DraftAngularDefect,
\end{equation}%
}
\newcommand{\CurvatureRegularDefectDisplay}{%
\begin{equation}
\DraftRegularDefect.
\end{equation}%
}
\newcommand{\CurvatureTotalDefectDisplay}{%
\begin{equation}
\DraftTotalDefect.
\label{eq:octahedral-total-defect}
\end{equation}%
}
\newcommand{\CurvatureInvariantDefectDisplay}{%
\begin{equation}
\DraftInvariantDefect.
\label{eq:defect}
\end{equation}%
}

% Energy

\newcommand{\EnergyTotalMassRelationDisplay}{%
\begin{equation}
\DraftTotalMassRelation.
\label{eq:mass}
\end{equation}%
}


% Spectrum


% CouplingMap
\newcommand{\CouplingMapPolynomialCouplingDisplay}{%
\begin{equation}
\PolynomialCouplingEquation.
\label{eq:polynomial}
\end{equation}%
}
\newcommand{\CouplingMapCouplingDifferentialDisplay}{%
\begin{equation}
\CouplingDifferentialEquation.
\label{eq:differential}
\end{equation}%
}
\newcommand{\CouplingMapCouplingSecondDifferentialDisplay}{%
\begin{equation}
\begin{gathered}
(D^2\Lambda(X)[\delta_1X,\delta_2X])_p
=\ExplainTerm{Dw_p(X)[\delta_1X](\delta_2X)_p}{first weight change}\\
+\ExplainTerm{Dw_p(X)[\delta_2X](\delta_1X)_p}{second weight change}\\
+\ExplainTerm{X_p\sum_{j\ne p}\bigl[\langle(\delta_1X)_p,(\delta_2X)_j\rangle_{\mathrm F}+\langle(\delta_2X)_p,(\delta_1X)_j\rangle_{\mathrm F}\bigr]}{second variation of the pairing}
\end{gathered}
\label{eq:second-differential}
\end{equation}%
}
\newcommand{\CouplingMapFullChainRuleDisplay}{%
\begin{equation}
\FullChainRuleEquation.
\label{eq:chain-rule}
\end{equation}%
}
\newcommand{\CouplingMapBothDifferentialVariationsDisplay}{%
\begin{equation}
\BothDifferentialVariationsEquation
\label{eq:both-differential-variations}
\end{equation}%
}
\newcommand{\CouplingMapPropagatedCompatibilityDisplay}{%
\begin{equation}
\PropagatedCompatibilityEquation
\label{eq:propagated-compatibility}
\end{equation}%
}
\newcommand{\CouplingMapContactDerivativeDisplay}{%
\begin{equation}
\ContactDerivativeEquation.
\label{eq:contact-derivative}
\end{equation}%
}

% Hilbert-space symbol used in the current Spectrum section.


% Full-source Spectrum repair. All formulas remain centralized here.
\newcommand{\SpectrumGraphMap}{\mathsf S}
\newcommand{\SpectrumQuotientMetric}{h_q}
\newcommand{\SpectrumGramMatrix}{\mathsf G}
\newcommand{\SpectrumNormalizer}{\mathsf N}

\newcommand{\SpectrumQuotientConnection}{A^{(q)}}
\newcommand{\SpectrumConnectionDifference}{A=A'+\Delta A}

\newcommand{\SpectrumActualStateMapDisplay}{%
\begin{equation}
\mathsf S^\dagger h_q\mathsf S=\mathsf G,\qquad
\mathsf N^\dagger\mathsf G\mathsf N=I,\qquad
U=h_q^{1/2}\mathsf S\mathsf N,\qquad U^\dagger U=I.
\label{eq:actual-state-map}
\end{equation}%
}
\newcommand{\SpectrumActualConnectionDisplay}{%
\begin{equation}
A'=U^{-1}A^{(q)}U+U^{-1}dU,\qquad
F'=U^{-1}F^{(q)}U,
\end{equation}%
}
\newcommand{\SpectrumActualCurvatureDisplay}{%
\begin{equation}
F_A=F'+D_{A'}(\Delta A)+\Delta A\wedge\Delta A.
\label{eq:actual-curvature-map}
\end{equation}%
}
\newcommand{\SpectrumFullCurvatureDisplay}{%
\begin{equation}
\begin{aligned}
\Omega_g&=\begin{pmatrix}\Omega_{11}&\Omega_{12}\\\Omega_{21}&0\end{pmatrix},&
\mathscr F&=\begin{pmatrix}F_A+\Delta F&(\mathscr F)_{12}\\(\mathscr F)_{21}&(\mathscr F)_{22}\end{pmatrix},\\
(\mathscr F)_{12}&=D_{\Omega_{11}}\Omega_{12},&(\mathscr F)_{21}&=D_{\Omega_{11}}\Omega_{21},\qquad (\mathscr F)_{22}=\Omega_{21}\wedge \Omega_{12}.
\end{aligned}
\end{equation}%
}
\newcommand{\SpectrumMixedCorrectionDisplay}{%
\begin{equation}
\Theta=A-\Omega_{11},\qquad \Delta F=\Omega_{12}\wedge \Omega_{21}-(D_{\Omega_{11}}\Theta+\Theta\wedge\Theta).
\end{equation}%
}

\newcommand{\SpectrumMetricConnection}{\Gamma(g)}

% Actual v5/v6 curvature-energy and sourced EYM chain.
\newcommand{\NativeCoefficientMetric}{\mathsf g}
\newcommand{\NativeSolder}{e}
\newcommand{\NativePairIndices}{A=(\mu,\nu),\ B=(\rho,\sigma)}
\newcommand{\NativeProjectors}{N^2=N,\ P=I-N}
\newcommand{\NativeRicciNotation}{\operatorname{Scal}_g=\operatorname{tr}_g\operatorname{Ric}}
\newcommand{\NativeWeylTensor}{\mathcal C}
\newcommand{\NativeWorkPi}{\Pi}
\newcommand{\NativeStressTrace}{\operatorname{tr}_g T=-\operatorname{Scal}_g/{\kappa}}
\newcommand{\NativeDegree}{\nu>0}
\newcommand{\NativeHeight}{R^2}
\newcommand{\NativeEuler}{\mathcal E}
\newcommand{\NativeEulerNormal}{n=\mathcal E/\sqrt{-g(\mathcal E,\mathcal E)}}
\newcommand{\NativeObserver}{u}
\newcommand{\NativeUnitNormal}{n}
\newcommand{\NativeNormalRicciZero}{\operatorname{Ric}(n,n)=0}
\newcommand{\NativeSpatialMetric}{h=g|_{n^\perp}}
\newcommand{\NativeSpatialTracefreeRicci}{\operatorname{Ric}_{0,ij}=\operatorname{Ric}_{ij}-(\operatorname{Scal}_g/3)h_{ij}}
\newcommand{\NativeRicciEigenvalues}{r_1,r_2,r_3}
\newcommand{\NativeSectionalValues}{k_{12},k_{13},k_{23}}
\newcommand{\NativeObservationVector}{\xi}
\newcommand{\NativeGeometryDisplay}{%
\begin{equation}
g=e^T\mathsf g e,\qquad \Gamma=e^{-1}(\Omega_g e+de).
\end{equation}%
}
\newcommand{\NativeCurvatureDisplay}{%
\begin{equation}
\mathscr F=d\Omega_g+\Omega_g\wedge\Omega_g,
\qquad \mathcal R=e^{-1}\mathscr F e.
\end{equation}%
}
\newcommand{\NativeEnergyDefinitionDisplay}{%
\begin{equation}
\begin{aligned}
(g^{-1}\wedge g^{-1})_{AB}&=g^{\mu\rho}g^{\nu\sigma}-g^{\mu\sigma}g^{\nu\rho},\\
\mathcal K&=-\sum_{A,B}(g^{-1}\wedge g^{-1})_{AB}\operatorname{tr}(\mathscr F_A\mathscr F_B).
\end{aligned}
\label{eq:native-energy}
\end{equation}%
}
\newcommand{\NativeRiemannEnergyDisplay}{%
\begin{equation}
\mathcal K=-\tfrac12\operatorname{tr}(\mathscr F_{\mu\nu}\mathscr F^{\mu\nu})=\tfrac12R_{abcd}R^{abcd}.
\end{equation}%
}
\newcommand{\NativePacketsDisplay}{%
\begin{equation}
\mathscr F_N=N\mathscr FN,\qquad
\mathscr F_C=N\mathscr FP+P\mathscr FN,\qquad
\mathscr F_T=P\mathscr FP.
\end{equation}%
}
\newcommand{\NativeEnergyAllocationDisplay}{%
\begin{equation}
\begin{gathered}
\mathcal K=\ExplainTerm{\mathcal K_N}{normal}+\ExplainTerm{\mathcal K_C}{mixed}+\ExplainTerm{\mathcal K_T}{tangent},\\
\mathcal K=(\mathcal K_N+\mathcal K_C)+(\mathcal K_C+\mathcal K_T)-\ExplainTerm{\mathcal K_C}{remove the duplicate mixed term}
\end{gathered}
\label{eq:native-allocation}
\end{equation}%
}
\newcommand{\NativeWeylEnergyDisplay}{%
\begin{equation}
\mathcal K=\tfrac12\langle\mathcal C,\mathcal C\rangle_g
+\langle \operatorname{Ric},\operatorname{Ric}\rangle_g-\tfrac16\operatorname{Scal}_g^2.
\label{eq:native-weyl-energy}
\end{equation}%
}
\newcommand{\NativeWorkCoefficientsDisplay}{%
\begin{equation}
\Pi=2\operatorname{Ric}-\tfrac13\operatorname{Scal}_g\,g,\qquad
\Pi_g=-2\operatorname{Ric}g^{-1}\operatorname{Ric}+\tfrac13\operatorname{Scal}_g \operatorname{Ric}.
\end{equation}%
}
\newcommand{\NativeFullWorkDisplay}{%
\begin{equation}
\delta \mathcal K=\delta\!\left(\tfrac12\langle\mathcal C,\mathcal C\rangle_g\right)+
\langle\Pi_g,\delta g\rangle_g+\langle\Pi,\delta \operatorname{Ric}\rangle_g.
\end{equation}%
}
\newcommand{\NativeStressExtractionDisplay}{%
\begin{equation}
T=\frac1{2{\kappa}}(\Pi-g\operatorname{tr}_g\Pi),\qquad
\operatorname{tr}_g\Pi=\tfrac23\operatorname{Scal}_g.
\end{equation}%
}
\newcommand{\NativeEinsteinDisplay}{%
\begin{equation}
{\kappa}T_{\mu\nu}=\operatorname{Ric}_{\mu\nu}-\tfrac12g_{\mu\nu}\operatorname{Scal}_g
=G_{\mu\nu}.
\label{eq:stress}
\end{equation}%
}
\newcommand{\NativeCurrentSourceDisplay}{%
\begin{equation}
\begin{aligned}
\operatorname{Ric}&={\kappa}\bigl(T-\tfrac12g\operatorname{tr}_gT\bigr),\\
J_{ab\nu}&=\nabla_a(\operatorname{Ric})_{b\nu}-\nabla_b(\operatorname{Ric})_{a\nu},
&j_\nu&=e\bigl(g^{ac}J_{cb\nu}\bigr)e^{-1}.
\end{aligned}
\end{equation}%
}
\newcommand{\NativeSourcedEquationsDisplay}{%
\begin{equation}
D^\mu\mathscr F_{\mu\nu}=j_\nu,\qquad\nabla^\mu T_{\mu\nu}=0.
\label{eq:native-sourced-eym}
\end{equation}%
}
\newcommand{\NativeEnergyStressDisplay}{%
\begin{equation}
\mathcal K=\tfrac12\langle\mathcal C,\mathcal C\rangle_g
+{\kappa}^2\left[\langle T,T\rangle_g-\tfrac16(\operatorname{tr}_gT)^2\right].
\label{eq:native-energy-stress}
\end{equation}%
}
\newcommand{\NativeTraceReversalNormDisplay}{%
\begin{equation}
\begin{aligned}
{\kappa}^2\langle T,T\rangle_g
&=\langle \operatorname{Ric},\operatorname{Ric}\rangle_g
-(\operatorname{Scal}_g)^2+(\operatorname{Scal}_g)^2\\
&=\langle \operatorname{Ric},\operatorname{Ric}\rangle_g.
\end{aligned}
\end{equation}%
}
\newcommand{\NativeHomothetyDisplay}{%
\begin{equation}
\mathcal L_{\mathcal E}g=4\nu g,\qquad
\mathcal E^\flat=-\nu\,d(R^2),\qquad
g(\mathcal E,\mathcal E)=-4\nu^2R^2<0.
\end{equation}%
}
\newcommand{\NativeEulerCurvatureDisplay}{%
\begin{equation}
\nabla_X\mathcal E=2\nu X,\qquad
\mathcal R(X,Y)\mathcal E
=2\nu(\nabla_XY-\nabla_YX-[X,Y])=0.
\end{equation}%
}
\newcommand{\NativeSpatialEnergyDisplay}{%
\begin{equation}
\mathcal K=2\sum_{i<j}\sum_{a<b}(R_{abij})^2\ge0,
\qquad \mathcal K=0\ \Longleftrightarrow\ \mathcal R=0.
\label{eq:native-spatial-positivity}
\end{equation}%
}
\newcommand{\NativeObservedEnergyDisplay}{%
\begin{equation}
E_\Sigma[u]=\int_\Sigma T_{\mu\nu}u^\mu n^\nu\,d\Sigma.
\label{eq:native-observed-energy}
\end{equation}%
}
\newcommand{\NativeObservedDensityDisplay}{%
\begin{equation}
\rho=T(n,n)=\frac1{\kappa}\left[\operatorname{Ric}(n,n)+\tfrac12\operatorname{Scal}_g\right].
\end{equation}%
}
\newcommand{\NativeSpatialDensityDisplay}{%
\begin{equation}
\rho=\frac{\operatorname{Scal}_g}{2{\kappa}}.
\label{eq:native-spatial-density}
\end{equation}%
}
\newcommand{\NativeDensityIdentityDisplay}{%
\begin{equation}
\mathcal K=2|\operatorname{Ric}_0|_h^2+\tfrac16\operatorname{Scal}_g^2=2|\operatorname{Ric}_0|_h^2+\tfrac23{\kappa}^2\rho^2.
\label{eq:native-density-identity}
\end{equation}%
}
\newcommand{\NativeRicciSectionalDisplay}{%
\begin{equation}
r_1=k_{12}+k_{13},\qquad r_2=k_{12}+k_{23},\qquad r_3=k_{13}+k_{23}.
\end{equation}%
}
\newcommand{\NativeSectionalSolutionDisplay}{%
\begin{equation}
k_{ij}=\tfrac12(r_i+r_j-r_k)=\tfrac12(\operatorname{Scal}_g-2r_k),\qquad\{i,j,k\}=\{1,2,3\}.
\end{equation}%
}
\newcommand{\NativeSpatialRicciEnergyDisplay}{%
\begin{equation}
\mathcal K=2\sum_{i<j}k_{ij}^2=2\sum_i r_i^2-\tfrac12\operatorname{Scal}_g^2.
\end{equation}%
}
\newcommand{\NativeChargeBalanceDisplay}{%
\begin{equation}
\nabla_\mu(T^{\mu\nu}\xi_\nu)=T^{\mu\nu}\nabla_{(\mu}\xi_{\nu)}.
\label{eq:native-charge-balance}
\end{equation}%
}


% Field definition preceding Phase.
\newcommand{\FieldVectorSpace}{W}
\newcommand{\FieldVectorPairing}{\langle\,\cdot\,,\,\cdot\,\rangle_W}
\newcommand{\FieldConfigurationDomain}{\mathcal U}
\newcommand{\FieldComponentSpace}{\mathbb V}
\newcommand{\FieldLabelSet}{\mathcal I}
\newcommand{\FieldSelfAdjointDisplay}{%
\begin{equation}
\underbrace{\langle Av,w\rangle}_{\text{act on }v\text{, then compare with }w}
=\underbrace{\langle v,Aw\rangle}_{\text{act on }w\text{, then compare with }v}
\qquad\text{for all }v,w.
\label{eq:field-self-adjoint}
\end{equation}%
}
\newcommand{\FieldMatrixPairingDisplay}{%
\begin{equation}
\langle A,B\rangle_{\mathrm F}=\operatorname{tr}(A^TB)=\operatorname{tr}(AB)=\sum_{a,b}A_{ab}B_{ab},\qquad A,B\in\mathbb V.
\label{eq:field-matrix-pairing}
\end{equation}%
}
\newcommand{\FieldDefinitionDisplay}{%
\begin{equation}
X\colon\mathcal U\longrightarrow {\mathcal V}=\mathbb V^{\mathcal I},\qquad
X(x)=\bigl(X_p(x)\bigr)_{p\in\mathcal I},
\label{eq:field-definition}
\end{equation}%
}
\newcommand{\FieldProductMetricDisplay}{%
\begin{equation}
\langle \delta_1X,\delta_2X\rangle_{\mathcal V}=\sum_{p\in\mathcal I}\langle (\delta_1X)_p,(\delta_2X)_p\rangle_{\mathrm F},\qquad \delta_1X,\delta_2X\in {\mathcal V}.
\label{eq:field-product-metric}
\end{equation}%
}

% Worked vector, matrix and trace comparisons in Field.
\newcommand{\FieldExampleVectorA}{a}
\newcommand{\FieldExampleVectorB}{b}
\newcommand{\FieldExampleMatrixA}{A}
\newcommand{\FieldExampleMatrixB}{B}
\newcommand{\FieldExampleMatrixTranspose}{A^T}
\newcommand{\FieldExampleMatrixSize}{2\times2}
\newcommand{\FieldTraceNotation}{\operatorname{tr}}
\newcommand{\FieldSymmetryIdentity}{A^T=A}
\newcommand{\FieldExampleVectorsDisplay}{%
\begin{equation}
a=(1,0),\qquad b=\left(\frac35,\frac45\right).
\label{eq:field-example-vectors}
\end{equation}%
}
\newcommand{\FieldExampleLengthsDisplay}{%
\begin{equation}
\langle a,a\rangle=1^2+0^2=1,\qquad
\langle b,b\rangle=\left(\frac35\right)^2+\left(\frac45\right)^2=1.
\end{equation}%
}
\newcommand{\FieldExampleAlignmentDisplay}{%
\begin{equation}
\langle a,b\rangle=1\cdot\frac35+0\cdot\frac45=\frac35.
\end{equation}%
}
\newcommand{\FieldExampleProjectionDisplay}{%
\begin{equation}
\underbrace{\left(\frac35,\frac45\right)}_{b}
=\underbrace{\left(\frac35,0\right)}_{\text{parallel to }a}
+\underbrace{\left(0,\frac45\right)}_{\text{perpendicular to }a}.
\end{equation}%
}
\newcommand{\FieldExampleMatricesDisplay}{%
\begin{equation}
A=\begin{pmatrix}a_{11}&a_{12}\\a_{21}&a_{22}\end{pmatrix},\qquad
B=\begin{pmatrix}b_{11}&b_{12}\\b_{21}&b_{22}\end{pmatrix},
\end{equation}%
}
\newcommand{\FieldExampleMatrixComparisonDisplay}{%
\begin{equation}
\langle A,B\rangle=a_{11}b_{11}+a_{12}b_{12}+a_{21}b_{21}+a_{22}b_{22}.
\end{equation}%
}
\newcommand{\FieldExampleMatrixLengthDisplay}{%
\begin{equation}
\|A\|^2=\langle A,A\rangle=a_{11}^{\,2}+a_{12}^{\,2}+a_{21}^{\,2}+a_{22}^{\,2}.
\end{equation}%
}
\newcommand{\FieldExampleTransposeDisplay}{%
\begin{equation}
A=\begin{pmatrix}a_{11}&a_{12}\\a_{21}&a_{22}\end{pmatrix}
\quad\longrightarrow\quad
A^T=\begin{pmatrix}a_{11}&a_{21}\\a_{12}&a_{22}\end{pmatrix}.
\end{equation}%
}
\newcommand{\FieldExampleDiagonalDisplay}{%
\begin{equation}
\begin{aligned}
(A^TB)_{11}&=a_{11}b_{11}+a_{21}b_{21},\\
(A^TB)_{22}&=a_{12}b_{12}+a_{22}b_{22}.
\end{aligned}
\end{equation}%
}
\newcommand{\FieldExampleTraceDisplay}{%
\begin{equation}
\begin{aligned}
\operatorname{tr}(A^TB)&=(A^TB)_{11}+(A^TB)_{22}\\
&=a_{11}b_{11}+a_{12}b_{12}+a_{21}b_{21}+a_{22}b_{22}\\
&=\langle A,B\rangle.
\end{aligned}
\label{eq:field-example-trace}
\end{equation}%
}

% Worked self-adjoint map in Field.
\newcommand{\FieldSelfAdjointExampleDisplay}{%
\begin{equation}
A=\begin{pmatrix}2&1\\1&3\end{pmatrix},\qquad
v=\begin{pmatrix}1\\0\end{pmatrix},\qquad
w=\begin{pmatrix}0\\1\end{pmatrix}.
\label{eq:field-self-adjoint-example}
\end{equation}%
}
\newcommand{\FieldSelfAdjointActionsDisplay}{%
\begin{equation}
Av=\begin{pmatrix}2\\1\end{pmatrix},\qquad
Aw=\begin{pmatrix}1\\3\end{pmatrix}.
\end{equation}%
}
\newcommand{\FieldSelfAdjointComparisonsDisplay}{%
\begin{equation}
\langle Av,w\rangle=2\cdot0+1\cdot1=1,\qquad
\langle v,Aw\rangle=1\cdot1+0\cdot3=1.
\end{equation}%
}
\newcommand{\FieldSelfAdjointIdentityDisplay}{%
\begin{equation}
\langle Av,w\rangle=(Av)^Tw=v^TA^Tw=v^TAw=\langle v,Aw\rangle.
\label{eq:field-self-adjoint-identity}
\end{equation}%
}

% Worked product-space metric in Field.
\newcommand{\FieldFirstChange}{(\delta_1X)_1}
\newcommand{\FieldSecondChange}{(\delta_1X)_2}
\newcommand{\FieldCompleteChange}{\delta_1X}
\newcommand{\FieldOtherChangePair}{\delta_2X=((\delta_2X)_1,(\delta_2X)_2)}
\newcommand{\FieldTwoComponentSpaceDisplay}{%
\begin{equation}
{\mathcal V}=\mathbb V\times\mathbb V.
\label{eq:field-two-component-space}
\end{equation}%
}
\newcommand{\FieldTwoComponentValueDisplay}{%
\begin{equation}
X=(X_1,X_2),\qquad X_1,X_2\in\mathbb V.
\end{equation}%
}
\newcommand{\FieldTwoComponentChangeDisplay}{%
\begin{equation}
\delta_1X=((\delta_1X)_1,(\delta_1X)_2).
\end{equation}%
}
\newcommand{\FieldTwoComponentMetricDisplay}{%
\begin{equation}
\langle \delta_1X,\delta_2X\rangle_{\mathcal V}=\langle (\delta_1X)_1,(\delta_2X)_1\rangle_{\mathrm F}+\langle (\delta_1X)_2,(\delta_2X)_2\rangle_{\mathrm F}.
\label{eq:field-two-component-metric}
\end{equation}%
}
\newcommand{\FieldProductFirstExampleDisplay}{%
\begin{equation}
(\delta_1X)_1=\begin{pmatrix}1&0\\0&0\end{pmatrix},\qquad
(\delta_1X)_2=\begin{pmatrix}0&0\\0&2\end{pmatrix},
\end{equation}%
}
\newcommand{\FieldProductSecondExampleDisplay}{%
\begin{equation}
(\delta_2X)_1=\begin{pmatrix}3&0\\0&0\end{pmatrix},\qquad
(\delta_2X)_2=\begin{pmatrix}0&0\\0&4\end{pmatrix}.
\end{equation}%
}
\newcommand{\FieldProductComparisonsDisplay}{%
\begin{equation}
\langle (\delta_1X)_1,(\delta_2X)_1\rangle_{\mathrm F}=1\cdot3=3,\qquad
\langle (\delta_1X)_2,(\delta_2X)_2\rangle_{\mathrm F}=2\cdot4=8.
\end{equation}%
}
\newcommand{\FieldProductTotalDisplay}{%
\begin{equation}
\langle \delta_1X,\delta_2X\rangle_{\mathcal V}=3+8=11.
\end{equation}%
}
\newcommand{\FieldProductLengthDisplay}{%
\begin{equation}
\|\delta_1X\|_{\mathcal V}^2=\langle (\delta_1X)_1,(\delta_1X)_1\rangle_{\mathrm F}+\langle (\delta_1X)_2,(\delta_1X)_2\rangle_{\mathrm F}=1^2+2^2=5.
\label{eq:field-product-example-length}
\end{equation}%
}

% Worked radial/transverse plane and U(1) rotations in Field.
\newcommand{\FieldRadialUnit}{{e_r}}
\newcommand{\FieldTransverseUnit}{{e_\perp}}
\newcommand{\FieldRadialProjectionValue}{\langle {e_r},b\rangle=\frac35}
\newcommand{\FieldTransverseLength}{\frac45}
\newcommand{\FieldPlaneComplexCoordinate}{z}
\newcommand{\FieldImaginaryUnit}{i}
\newcommand{\FieldPlaneRealCoordinates}{(x,y)}
\newcommand{\FieldPlaneRotatedCoordinates}{(-y,x)}
\newcommand{\FieldPlaneComplexStructure}{J}
\newcommand{\FieldPlaneIdentity}{I}
\newcommand{\FieldPlaneAngle}{\theta}
\newcommand{\FieldRadialVectorsDisplay}{%
\begin{equation}
a=(1,0),\qquad b=\left(\frac35,\frac45\right).
\end{equation}%
}
\newcommand{\FieldRadialUnitDisplay}{%
\begin{equation}
{e_r}=a=(1,0).
\end{equation}%
}
\newcommand{\FieldTransverseSubtractionDisplay}{%
\begin{equation}
b_\perp=b-\langle {e_r},b\rangle {e_r}
=\left(\frac35,\frac45\right)-\left(\frac35,0\right)
=\left(0,\frac45\right).
\label{eq:field-transverse-example}
\end{equation}%
}
\newcommand{\FieldTransverseUnitDisplay}{%
\begin{equation}
{e_\perp}=\frac{b_\perp}{\|b_\perp\|}=(0,1).
\end{equation}%
}
\newcommand{\FieldPlaneOrthonormalityDisplay}{%
\begin{equation}
\langle {e_r},{e_r}\rangle=\langle {e_\perp},{e_\perp}\rangle=1,\qquad \langle {e_r},{e_\perp}\rangle=0.
\end{equation}%
}
\newcommand{\FieldPlaneVectorDisplay}{%
\begin{equation}
v=x{e_r}+y{e_\perp}.
\end{equation}%
}
\newcommand{\FieldPlaneComplexCoordinateDisplay}{%
\begin{equation}
z=x+iy,\qquad i^2=-1.
\end{equation}%
}
\newcommand{\FieldPlaneImaginaryRotationDisplay}{%
\begin{equation}
iz=i(x+iy)=-y+ix.
\end{equation}%
}
\newcommand{\FieldPlaneQuarterTurnDisplay}{%
\begin{equation}
{e_r}\longmapsto {e_\perp},\qquad {e_\perp}\longmapsto -{e_r}.
\end{equation}%
}
\newcommand{\FieldPlaneComplexStructureDisplay}{%
\begin{equation}
J{e_r}={e_\perp},\qquad J{e_\perp}=-{e_r},\qquad J^2=-I,
\end{equation}%
}
\newcommand{\FieldPlanePhaseFactorDisplay}{%
\begin{equation}
e^{i\theta}=\cos\theta+i\sin\theta.
\end{equation}%
}
\newcommand{\FieldPlaneRotationCoordinatesDisplay}{%
\begin{equation}
e^{i\theta}(x+iy)
=(x\cos\theta-y\sin\theta)+i(x\sin\theta+y\cos\theta),
\end{equation}%
}
\newcommand{\FieldPlanePhaseLengthDisplay}{%
\begin{equation}
|e^{i\theta}|^2=\cos^2\theta+\sin^2\theta=1,
\end{equation}%
}
\newcommand{\FieldPlanePhaseGroupDisplay}{%
\begin{equation}
\mathrm U(1)=\{e^{i\theta}:\theta\in\mathbb R\}.
\label{eq:field-phase-group-example}
\end{equation}%
}
\newcommand{\FieldPlanePhaseCompositionDisplay}{%
\begin{equation}
e^{i\alpha}e^{i\beta}=e^{i(\alpha+\beta)}.
\end{equation}%
}

% Worked induced U(1) connection and local gauge transformation in Field.
\newcommand{\FieldPathParameter}{s}
\newcommand{\FieldMovingRadial}{{e_r}(s)}
\newcommand{\FieldMovingTransverse}{{e_\perp}(s)}
\newcommand{\FieldRadialDerivative}{d e_r/ds}
\newcommand{\FieldTransverseDerivative}{d e_\perp/ds}
\newcommand{\FieldNegativeRadial}{-{e_r}}
\newcommand{\FieldConnectionComplexValue}{z=x+iy}
\newcommand{\FieldGaugeAngle}{\chi(s)}
\newcommand{\FieldConnectionVector}{v}
\newcommand{\FieldGaugeUnrotation}{\chi=-s}
\newcommand{\FieldConfigurationCoordinates}{x^\mu}
\newcommand{\FieldConnectionCoefficients}{a_\mu}
\newcommand{\FieldConnectionRateDisplay}{%
\begin{equation}
 a_s=\left\langle e_\perp(s),\frac{d e_r}{ds}\right\rangle.
\label{eq:field-connection-rate}
\end{equation}%
}
\newcommand{\FieldConnectionProjectedDisplay}{%
\begin{equation}
\left(\frac{d e_r}{ds}\right)_{\parallel}=a_s {e_\perp},\qquad
\left(\frac{d e_\perp}{ds}\right)_{\parallel}=-a_s {e_r}.
\end{equation}%
}
\newcommand{\FieldConnectionRotatingFrameDisplay}{%
\begin{equation}
{e_r}(s)=(\cos s,\sin s),\qquad {e_\perp}(s)=(-\sin s,\cos s).
\end{equation}%
}
\newcommand{\FieldConnectionRotatingRateDisplay}{%
\begin{equation}
\frac{d e_r}{ds}={e_\perp},\qquad\frac{d e_\perp}{ds}=-{e_r},\qquad a_s=\langle {e_\perp},{e_\perp}\rangle=1.
\end{equation}%
}
\newcommand{\FieldConnectionVectorDisplay}{%
\begin{equation}
v(s)=x(s){e_r}(s)+y(s){e_\perp}(s).
\end{equation}%
}
\newcommand{\FieldConnectionVectorDerivativeDisplay}{%
\begin{equation}
\left(\frac{dv}{ds}\right)_{\parallel}
=\left(\frac{dx}{ds}-a_s y\right){e_r}
+\left(\frac{dy}{ds}+a_s x\right){e_\perp}.
\end{equation}%
}
\newcommand{\FieldConnectionCovariantDerivativeDisplay}{%
\begin{equation}
D_s z=\frac{dz}{ds}+ia_s z.
\end{equation}%
}
\newcommand{\FieldConnectionFixedVectorDisplay}{%
\begin{equation}
v=(1,0).
\end{equation}%
}
\newcommand{\FieldConnectionFixedCoordinatesDisplay}{%
\begin{equation}
x(s)=\langle {e_r}(s),v\rangle=\cos s,\qquad
y(s)=\langle {e_\perp}(s),v\rangle=-\sin s,
\end{equation}%
}
\newcommand{\FieldConnectionFixedComplexDisplay}{%
\begin{equation}
z(s)=e^{-is}.
\end{equation}%
}
\newcommand{\FieldConnectionCancellationDisplay}{%
\begin{equation}
D_sz=\frac{d}{ds}e^{-is}+ie^{-is}=-ie^{-is}+ie^{-is}=0.
\label{eq:field-connection-fixed-vector}
\end{equation}%
}
\newcommand{\FieldConnectionFrameChangeDisplay}{%
\begin{equation}
\widetilde {e_r}=\cos\chi\,{e_r}+\sin\chi\,{e_\perp},\qquad
\widetilde {e_\perp}=-\sin\chi\,{e_r}+\cos\chi\,{e_\perp}.
\end{equation}%
}
\newcommand{\FieldConnectionCoordinateChangeDisplay}{%
\begin{equation}
\widetilde z=e^{-i\chi}z.
\end{equation}%
}
\newcommand{\FieldConnectionPotentialChangeDisplay}{%
\begin{equation}
\widetilde a_s=a_s+\frac{d\chi}{ds}.
\end{equation}%
}
\newcommand{\FieldConnectionCovarianceDisplay}{%
\begin{equation}
\widetilde D_s\widetilde z=e^{-i\chi}D_sz.
\label{eq:field-connection-gauge-covariance}
\end{equation}%
}
\newcommand{\FieldConnectionCoordinateDirectionsDisplay}{%
\begin{equation}
a_\mu=\left\langle {e_\perp},\frac{\partial {e_r}}{\partial x^\mu}\right\rangle,\qquad
D_\mu z=\frac{\partial z}{\partial x^\mu}+ia_\mu z.
\end{equation}%
}

% Worked coupling, orthogonal planes and non-Abelian frame transformations.
\newcommand{\FieldCouplingLinearResponse}{{L_X}\delta X}
\newcommand{\FieldCouplingExampleAdjoint}{{L_X}^*=2I}
\newcommand{\FieldCouplingExampleNormalization}{\sqrt5}
\newcommand{\FieldCouplingUnitaryGroup}{\mathrm U(2)}
\newcommand{\FieldCouplingMixingU}{U_1}
\newcommand{\FieldCouplingMixingV}{U_2}
\newcommand{\FieldCouplingNoncommutation}{U_1U_2\ne U_2U_1}
\newcommand{\FieldCouplingResponseDisplay}{%
\begin{equation}Y=\Lambda(X).\end{equation}%
}
\newcommand{\FieldCouplingWeightsDisplay}{%
\begin{equation}
Y_p=w_p(X)X_p,\qquad w_p(X)=\sum_{j\ne p}\langle X_p,X_j\rangle_{\mathrm F}.
\end{equation}%
}
\newcommand{\FieldCouplingDerivativeDisplay}{%
\begin{equation}{L_X}=D\Lambda(X).\end{equation}%
}
\newcommand{\FieldCouplingVariationDisplay}{%
\begin{equation}
({L_X}\delta X)_p=w_p(X)(\delta X)_p+X_p\sum_{j\ne p}\left[\langle (\delta X)_p,X_j\rangle_{\mathrm F}+\langle X_p,(\delta X)_j\rangle_{\mathrm F}\right].
\label{eq:field-coupling-variation}
\end{equation}%
}
\newcommand{\FieldCouplingGraphDisplay}{%
\begin{equation}{\iota_{L_X}}(\delta X)=(\delta X,{L_X}\delta X).\end{equation}%
}
\newcommand{\FieldCouplingAdjointDisplay}{%
\begin{equation}
\langle {L_X}\delta X,w\rangle_{\mathcal V}=\langle \delta X,{L_X}^*w\rangle_{\mathcal V}\qquad\text{for every }\delta X,w.
\end{equation}%
}
\newcommand{\FieldCouplingNormalDisplay}{%
\begin{equation}{\iota_{L_X}^\perp}(w)=(-{L_X}^*w,w).\end{equation}%
}
\newcommand{\FieldCouplingOrthogonalityDisplay}{%
\begin{equation}
\begin{aligned}
\langle {\iota_{L_X}}(\delta X),{\iota_{L_X}^\perp}(w)\rangle_{{\mathcal V}\times {\mathcal V}}
&=\langle \delta X,-{L_X}^*w\rangle_{\mathcal V}+\langle {L_X}\delta X,w\rangle_{\mathcal V}\\
&=-\langle \delta X,{L_X}^*w\rangle_{\mathcal V}+\langle {L_X}\delta X,w\rangle_{\mathcal V}\\
&=0.
\end{aligned}
\label{eq:field-coupling-orthogonality}
\end{equation}%
}
\newcommand{\FieldCouplingGraphPlaneDisplay}{%
\begin{equation}{\iota_{L_X}}({e_r})=({e_r},{L_X}{e_r}),\qquad {\iota_{L_X}}({e_\perp})=({e_\perp},{L_X}{e_\perp}),\end{equation}%
}
\newcommand{\FieldCouplingNormalPlaneDisplay}{%
\begin{equation}{\iota_{L_X}^\perp}({e_r})=(-{L_X}^*{e_r},{e_r}),\qquad {\iota_{L_X}^\perp}({e_\perp})=(-{L_X}^*{e_\perp},{e_\perp}).\end{equation}%
}
\newcommand{\FieldCouplingExampleDerivativeDisplay}{%
\begin{equation}{L_X}=2I,\qquad {e_r}=(1,0),\qquad {e_\perp}=(0,1),\end{equation}%
}
\newcommand{\FieldCouplingExampleLiftsDisplay}{%
\begin{equation}
\begin{aligned}
{\iota_{L_X}}({e_r})&=(1,0;\,2,0),&{\iota_{L_X}}({e_\perp})&=(0,1;\,0,2),\\
{\iota_{L_X}^\perp}({e_r})&=(-2,0;\,1,0),&{\iota_{L_X}^\perp}({e_\perp})&=(0,-2;\,0,1).
\end{aligned}
\end{equation}%
}
\newcommand{\FieldCouplingExampleCancellationDisplay}{%
\begin{equation}
\langle {\iota_{L_X}}({e_r}),{\iota_{L_X}^\perp}({e_r})\rangle=1(-2)+0(0)+2(1)+0(0)=0.
\end{equation}%
}
\newcommand{\FieldCouplingComplexPairDisplay}{%
\begin{equation}\mathbf z=\begin{pmatrix}z_1\\z_2\end{pmatrix},\end{equation}%
}
\newcommand{\FieldCouplingComplexNormDisplay}{%
\begin{equation}\|\mathbf z\|^2=|z_1|^2+|z_2|^2.\end{equation}%
}
\newcommand{\FieldCouplingMixingMatricesDisplay}{%
\begin{equation}
{U_1}=\begin{pmatrix}i&0\\0&-i\end{pmatrix},\qquad
{U_2}=\frac1{\sqrt2}\begin{pmatrix}1&1\\-1&1\end{pmatrix}.
\end{equation}%
}
\newcommand{\FieldCouplingInitialCoordinateDisplay}{%
\begin{equation}\mathbf z_0=\begin{pmatrix}1\\0\end{pmatrix},\end{equation}%
}
\newcommand{\FieldCouplingFirstOrderDisplay}{%
\begin{equation}{U_1}{U_2}\mathbf z_0=\frac1{\sqrt2}\begin{pmatrix}i\\i\end{pmatrix},\end{equation}%
}
\newcommand{\FieldCouplingSecondOrderDisplay}{%
\begin{equation}{U_2}{U_1}\mathbf z_0=\frac1{\sqrt2}\begin{pmatrix}i\\-i\end{pmatrix}.\label{eq:field-coupling-noncommuting-example}\end{equation}%
}

% Yang--Mills waves, collective mass and the physical mass-gap question.
\newcommand{\FieldYMPotential}{A_\mu}
\newcommand{\FieldYMCoordinateIndex}{\mu}
\newcommand{\FieldYMCoupling}{g_{\mathrm{YM}}}
\newcommand{\FieldYMNaturalUnits}{c=\hbar=1}
\newcommand{\FieldYMMassTerm}{\operatorname{tr}(A_\mu A^\mu)}
\newcommand{\FieldYMFrequency}{\omega}
\newcommand{\FieldYMWaveVector}{\mathbf k}
\newcommand{\FieldYMQuantumEnergy}{E=\hbar\omega}
\newcommand{\FieldYMQuantumMomentum}{\mathbf p=\hbar\mathbf k}
\newcommand{\FieldYMZeroMass}{m=0}
\newcommand{\FieldYMTotalEnergy}{E}
\newcommand{\FieldYMTotalMomentum}{\mathbf P}
\newcommand{\FieldYMCollectiveMass}{M}
\newcommand{\FieldYMExampleEnergy}{\varepsilon}
\newcommand{\FieldYMScalarLabel}{0^{++}}


\newcommand{\FieldYMStrengthDisplay}{%
\begin{equation}
F_{\mu\nu}=\partial_\mu A_\nu-\partial_\nu A_\mu-ig_{\mathrm{YM}}[A_\mu,A_\nu],
\label{eq:field-ym-strength}
\end{equation}%
}
\newcommand{\FieldYMCommutatorDisplay}{%
\begin{equation}[A_\mu,A_\nu]=A_\mu A_\nu-A_\nu A_\mu.\end{equation}%
}
\newcommand{\FieldYMLagrangianDisplay}{%
\begin{equation}\mathcal L_{\mathrm{YM}}=-\frac12\operatorname{tr}(F_{\mu\nu}F^{\mu\nu}).\end{equation}%
}
\newcommand{\FieldYMWaveDisplay}{%
\begin{equation}
\left(\frac1{c^2}\frac{\partial^2}{\partial t^2}-\nabla^2\right)\mathbf A_\perp=0.
\end{equation}%
}
\newcommand{\FieldYMDispersionDisplay}{%
\begin{equation}\omega^2=c^2|\mathbf k|^2.\end{equation}%
}
\newcommand{\FieldYMMasslessEnergyDisplay}{%
\begin{equation}E=c|\mathbf p|.\end{equation}%
}
\newcommand{\FieldYMRelativisticEnergyDisplay}{%
\begin{equation}E^2=c^2|\mathbf p|^2+m^2c^4.\end{equation}%
}
\newcommand{\FieldYMInvariantMassDisplay}{%
\begin{equation}M^2c^4=E^2-c^2|\mathbf P|^2.\end{equation}%
}
\newcommand{\FieldYMOppositeMomentaDisplay}{%
\begin{equation} E_1=E_2=\varepsilon,\qquad \mathbf p_1=\frac{\varepsilon}{c}\hat{\mathbf x},\qquad \mathbf p_2=-\frac{\varepsilon}{c}\hat{\mathbf x},\end{equation}%
}
\newcommand{\FieldYMCombinedEnergyDisplay}{%
\begin{equation} E=E_1+E_2=2\varepsilon,\qquad\mathbf P=0.\end{equation}%
}
\newcommand{\FieldYMPairMassDisplay}{%
\begin{equation} M=\frac{2\varepsilon}{c^2}.\label{eq:field-massless-pair-mass}\end{equation}%
}
\newcommand{\FieldYMGlueballMassDisplay}{%
\begin{equation}M_{0^{++}}c^2=1.73\pm0.05\pm0.08\ \mathrm{GeV}\end{equation}%
}


% Worked Phase applications, loop phases and gauge invariance.
\newcommand{\PhaseExampleRotationAngle}{\pi/4}
\newcommand{\PhaseHermitianPairing}{h}
\newcommand{\PhaseDifferentialSymbol}{d}
\newcommand{\PhaseOrientedAreaElement}{d x^1\wedge d x^2}
\newcommand{\PhaseSurfaceBoundary}{\partial\mathcal S}
\newcommand{\PhaseExampleLoopAngle}{\pi/2}
\newcommand{\PhaseConnectionConvention}{\nabla=d+ia}
\newcommand{\PhaseGaugeAngleSymbol}{\chi}
\newcommand{\PhaseMatrixExampleDisplay}{%
\begin{equation}
X_p=\begin{pmatrix}1&0\\0&0\end{pmatrix},\qquad
X_i=\begin{pmatrix}\frac35&0\\0&\frac45\end{pmatrix}.
\label{eq:phase-matrix-example}
\end{equation}%
}
\newcommand{\PhaseMatrixProjectionDisplay}{%
\begin{equation}
\begin{gathered}
\langle X_p,X_p\rangle_{\mathrm F}=1,\qquad\langle X_p,X_i\rangle_{\mathrm F}=\frac35,\qquad
X_{i\perp p}=\begin{pmatrix}0&0\\0&\frac45\end{pmatrix},\\
{e_r}=\begin{pmatrix}1&0\\0&0\end{pmatrix},\qquad
{e_\perp}=\begin{pmatrix}0&0\\0&1\end{pmatrix},\qquad\langle {e_r},{e_\perp}\rangle_{\mathrm F}=0.
\end{gathered}
\end{equation}%
}
\newcommand{\PhaseMatrixRotationDisplay}{%
\begin{equation}
\mathsf R\!\left(\frac\pi4\right){e_r}=\frac{{e_r}+{e_\perp}}{\sqrt2}
=\frac1{\sqrt2}\begin{pmatrix}1&0\\0&1\end{pmatrix},\qquad
\left\|\mathsf R\!\left(\frac\pi4\right){e_r}\right\|_{\mathrm F}^2=1.
\end{equation}%
}
\newcommand{\PhaseHermitianExampleDisplay}{%
\begin{equation}
\begin{gathered}
v=x{e_r}+y{e_\perp},\qquad z=x+iy,\qquad\langle v,v\rangle_{\mathrm F}=x^2+y^2=|z|^2,\\
h(z_1,z_2)=\overline{z_1}z_2.
\end{gathered}
\end{equation}%
}
\newcommand{\PhaseConnectionPathDisplay}{%
\begin{equation}
a_s=a\!\left(\frac{d}{ds}\right)=\left\langle e_\perp,\frac{d e_r}{ds}\right\rangle_{\mathrm F}.
\end{equation}%
}
\newcommand{\PhaseCurvatureCoordinatesDisplay}{%
\begin{equation}
a=a_1\,dx^1+a_2\,dx^2,\qquad
da=\left(\frac{\partial a_2}{\partial x^1}-\frac{\partial a_1}{\partial x^2}\right)dx^1\wedge dx^2.
\label{eq:phase-curvature-coordinates}
\end{equation}%
}
\newcommand{\PhaseLoopExampleDisplay}{%
\begin{equation}
\int_{\mathcal S} da=\frac\pi2\quad\Longrightarrow\quad
\exp\!\left(i\oint_{\partial \mathcal S}a\right)=e^{i\pi/2}=i.
\end{equation}%
}
\newcommand{\PhaseParallelTransportDisplay}{%
\begin{equation}
dz+iaz=0,\qquad
z_{\mathrm{end}}=\exp\!\left(-i\oint_{\partial\mathcal S}a\right)z_{\mathrm{start}}.
\label{eq:phase-parallel-transport-sign}
\end{equation}%
}
\newcommand{\PhaseGaugeExampleDisplay}{%
\begin{equation}
\chi(x^1,x^2)=x^1x^2,\qquad d\chi=x^2\,dx^1+x^1\,dx^2,\qquad d(d\chi)=(-1+1)\,dx^1\wedge dx^2=0.
\end{equation}%
}
\newcommand{\PhaseGaugeAngleDifferential}{d\chi}

% Worked Coupling examples and explanations.
\newcommand{\CouplingPlaneCoordinates}{\xi}
\newcommand{\CouplingAlpha}{\alpha}
\newcommand{\CouplingBeta}{\beta}
\newcommand{\CouplingSecondDerivative}{{L_Y}}
\newcommand{\CouplingMatrixInputExampleDisplay}{%
\begin{equation}
X_1=\begin{pmatrix}1&0\\0&0\end{pmatrix},\qquad
X_2=\begin{pmatrix}2&0\\0&1\end{pmatrix}.
\label{eq:coupling-matrix-input}
\end{equation}%
}
\newcommand{\CouplingMatrixResponseExampleDisplay}{%
\begin{equation}
\begin{gathered}
w_1=w_2=\langle X_1,X_2\rangle_{\mathrm F}=2,\\
Y_1=2X_1=\begin{pmatrix}2&0\\0&0\end{pmatrix},\qquad
Y_2=2X_2=\begin{pmatrix}4&0\\0&2\end{pmatrix}.
\end{gathered}
\end{equation}%
}
\newcommand{\CouplingMatrixVariationExampleDisplay}{%
\begin{equation}
\begin{gathered}
(\delta X)_1=\begin{pmatrix}0&0\\0&1\end{pmatrix},\qquad (\delta X)_2=0,\qquad
\delta w_1=\delta w_2=\langle (\delta X)_1,X_2\rangle_{\mathrm F}=1,\\
(L_X\delta X)_1=2(\delta X)_1+X_1=\begin{pmatrix}1&0\\0&2\end{pmatrix},\qquad
(L_X\delta X)_2=X_2=\begin{pmatrix}2&0\\0&1\end{pmatrix}.
\end{gathered}
\label{eq:coupling-matrix-variation}
\end{equation}%
}
\newcommand{\CouplingAdjointMeaningDisplay}{%
\begin{equation}
\langle {L_X}v,w\rangle_{\mathcal V}=\langle v,{L_X}^*w\rangle_{\mathcal V}\qquad(v,w\in {\mathcal V}).
\end{equation}%
}
\newcommand{\CouplingGramLengthsDisplay}{%
\begin{equation}
\begin{aligned}
\xi^T{\mathsf V_0}^*{\mathsf V_0}\xi&=\|{\mathsf V}\xi\|_{\mathcal V}^2+\|{L_X}{\mathsf V}\xi\|_{\mathcal V}^2>0,\\
\xi^T{\mathsf V_1}^*{\mathsf V_1}\xi&=\|{\mathsf V}\xi\|_{\mathcal V}^2+\|{L_X}^*{\mathsf V}\xi\|_{\mathcal V}^2>0,
\qquad\xi\in\mathbb R^2\setminus\{0\}.
\end{aligned}
\label{eq:coupling-gram-lengths}
\end{equation}%
}
\newcommand{\CouplingComplexVolumeDisplay}{%
\begin{equation}
\det[z\ w]=z_1w_2-z_2w_1,\qquad
\det[Uz\ Uw]=(\det U)\det[z\ w].
\end{equation}%
}
\newcommand{\CouplingUnitaryMixingExampleDisplay}{%
\begin{equation}
\alpha=\beta=\frac1{\sqrt2},\qquad
U=\frac1{\sqrt2}\begin{pmatrix}1&1\\-1&1\end{pmatrix},\qquad
U\begin{pmatrix}1\\0\end{pmatrix}=\frac1{\sqrt2}\begin{pmatrix}1\\-1\end{pmatrix}.
\end{equation}%
}
\newcommand{\CouplingSuccessiveResponseDisplay}{%
\begin{equation}
Y=\Lambda(X),\quad Z=\Lambda(Y),\quad {L_Y}=D\Lambda(Y),\qquad
\delta Y={L_X}\delta X,\quad\delta Z={L_Y}{L_X}\delta X.
\label{eq:coupling-successive-response}
\end{equation}%
}
\newcommand{\CouplingSecondMatrixResponseDisplay}{%
\begin{equation}
\langle Y_1,Y_2\rangle_{\mathrm F}=8,\qquad
Z_1=8Y_1=\begin{pmatrix}16&0\\0&0\end{pmatrix},\qquad
Z_2=8Y_2=\begin{pmatrix}32&0\\0&16\end{pmatrix}.
\end{equation}%
}

% Worked explanations for Algebra.
\newcommand{\AlgebraRawNormal}{{\mathsf W_2}}
\newcommand{\AlgebraFinalNormal}{{\mathsf V_2}}
\newcommand{\AlgebraResponseExampleDisplay}{%
\begin{equation}
\begin{aligned}
\delta w_1(Y)=\delta w_2(Y)
&=\langle(L_X\delta X)_1,Y_2\rangle_{\mathrm F}+\langle Y_1,(L_X\delta X)_2\rangle_{\mathrm F}\\
&=8+4=12,\\
(L_YL_X\delta X)_1&=8(L_X\delta X)_1+12Y_1=\begin{pmatrix}32&0\\0&16\end{pmatrix},\\
(L_YL_X\delta X)_2&=8(L_X\delta X)_2+12Y_2=\begin{pmatrix}64&0\\0&32\end{pmatrix}.
\end{aligned}
\label{eq:algebra-response-example}
\end{equation}%
}
\newcommand{\AlgebraNormalCancellationDisplay}{%
\begin{equation}
\begin{aligned}
\langle\iota{}v,{\iota_1^\perp}w\rangle
&=-\langle v,{L_X}^*w\rangle+\langle {L_X}v,w\rangle=0,\\
\langle\iota{}v,{\iota_2^\perp}w\rangle
&=-\langle {L_X}v,{L_Y}^*w\rangle+\langle {L_Y}{L_X}v,w\rangle=0.
\end{aligned}
\end{equation}%
}
\newcommand{\AlgebraProjectionDisplay}{%
\begin{equation}
({\mathsf V_1}^*{\mathsf V_1})\alpha={\mathsf V_1}^*{\mathsf W_2}\xi,\qquad \alpha=({\mathsf V_1}^*{\mathsf V_1})^{-1}{\mathsf V_1}^*{\mathsf W_2}\xi,\qquad
{\mathsf V_1}^*({\mathsf W_2}\xi-{\mathsf V_1}\alpha)=0.
\end{equation}%
}
\newcommand{\AlgebraMatrixUnitExampleDisplay}{%
\begin{equation}
E_{12}\begin{pmatrix}z_1\\z_2\\z_3\end{pmatrix}
=\begin{pmatrix}z_2\\0\\0\end{pmatrix},\qquad
E_{23}\begin{pmatrix}z_1\\z_2\\z_3\end{pmatrix}
=\begin{pmatrix}0\\z_3\\0\end{pmatrix},\qquad
E_{12}E_{23}=E_{13}.
\end{equation}%
}
\newcommand{\AlgebraGeneratorActionDisplay}{%
\begin{equation}
\mathsf A_{12}=\begin{pmatrix}0&1&0\\-1&0&0\\0&0&0\end{pmatrix},\qquad
\mathsf A_{12}\begin{pmatrix}z_1\\z_2\\z_3\end{pmatrix}
=\begin{pmatrix}z_2\\-z_1\\0\end{pmatrix},\qquad
\mathsf A_{12}^{\dagger}=-\mathsf A_{12}.
\end{equation}%
}
\newcommand{\AlgebraCommutatorMeaningDisplay}{%
\begin{equation}
[P,Q]=PQ-QP,\qquad [P,Q]z=P(Qz)-Q(Pz).
\end{equation}%
}
\newcommand{\AlgebraBracketProductsDisplay}{%
\begin{equation}
\mathsf A_{12}\mathsf A_{23}=E_{13},\qquad \mathsf A_{23}\mathsf A_{12}=E_{31},\qquad
[\mathsf A_{12},\mathsf A_{23}]=E_{13}-E_{31}=\mathsf A_{13}.
\end{equation}%
}
\newcommand{\AlgebraDimensionDisplay}{%
\begin{equation}
\operatorname{diag}(ia_1,ia_2,ia_3),\qquad
 a_3=-a_1-a_2,\qquad
\dim_{\mathbb R}\mathfrak{su}(3)=3\cdot2+(3-1)=8.
\end{equation}%
}
\newcommand{\AlgebraFrameSkewDisplay}{%
\begin{equation}
\mathsf E^*\mathsf E=I_6,\qquad
 d(\mathsf E^*\mathsf E)=(d\mathsf E)^*\mathsf E+\mathsf E^*d\mathsf E=0,\qquad
\Omega_{\mathbb R}^*=-\Omega_{\mathbb R}.
\end{equation}%
}
\newcommand{\AlgebraConnectionPropertiesDisplay}{%
\begin{equation}
A_JJ=JA_J,\qquad A^\dagger=-A,\qquad \operatorname{tr}_{\mathbb C}A=0.
\end{equation}%
}
\newcommand{\AlgebraWedgeExampleDisplay}{%
\begin{equation}
\begin{gathered}
(A_{12}\wedge A_{23})(\eta,\zeta)=A_{12}(\eta)A_{23}(\zeta)-A_{12}(\zeta)A_{23}(\eta),\\
A_{12}(\eta)=A_{23}(\zeta)=1,\quad A_{12}(\zeta)=A_{23}(\eta)=0
\quad\Longrightarrow\quad (A_{12}\wedge A_{23})(\eta,\zeta)=1.
\end{gathered}
\end{equation}%
}
\newcommand{\AlgebraDiagonalNormExampleDisplay}{%
\begin{equation}
\tfrac12\operatorname{tr}(\lambda_3^2)=\frac{1+1}{2}=1,\quad \tfrac12\operatorname{tr}(\lambda_8^2)=\frac{1+1+4}{6}=1,\quad \tfrac12\operatorname{tr}(\lambda_3\lambda_8)=\frac{1-1}{\sqrt{12}}=0.
\end{equation}%
}
\newcommand{\AlgebraDiagonalExpansionDisplay}{%
\begin{equation}
\operatorname{diag}(1,0,-1)=\frac{\lambda_3}{2}+\frac{\sqrt3\,\lambda_8}{2}=\operatorname{diag}\!\left(\tfrac12,-\tfrac12,0\right)+\operatorname{diag}\!\left(\tfrac12,\tfrac12,-1\right).
\end{equation}%
}

% Worked explanations for Curvature.
\newcommand{\CurvatureEdgeLength}{\sqrt2\,\ell}
\newcommand{\CurvatureEndpointExampleDisplay}{%
\begin{equation}
(\ell,0,0),\quad(-\ell,0,0),\quad(0,\ell,0),\quad(0,-\ell,0),\quad(0,0,\ell),\quad(0,0,-\ell).
\end{equation}%
}
\newcommand{\CurvatureFaceAverageDisplay}{%
\begin{equation}
\begin{gathered}
x=\alpha\ell e_1+\beta\ell e_2+\gamma\ell e_3,
\qquad\alpha,\beta,\gamma\ge0,\qquad\alpha+\beta+\gamma=1,\\
\alpha=\beta=\gamma=\tfrac13
\quad\Longrightarrow\quad x=\tfrac{\ell}{3}(e_1+e_2+e_3).
\end{gathered}
\end{equation}%
}
\newcommand{\CurvatureFacePairExampleDisplay}{%
\begin{equation}
\tau_{(+,+,+)}=\operatorname{conv}\{\ell e_1,\ell e_2,\ell e_3\},\qquad
\tau_{(+,+,-)}=\operatorname{conv}\{\ell e_1,\ell e_2,-\ell e_3\}.
\end{equation}%
}
\newcommand{\CurvatureEdgeCancellationExampleDisplay}{%
\begin{equation}
[\ell e_1,\ell e_2]+[\ell e_2,\ell e_1]=0.
\end{equation}%
}
\newcommand{\CurvatureCountingExampleDisplay}{%
\begin{equation}
N_{\mathrm{vertices}}=3\cdot2=6,\qquad
N_{\mathrm{edges}}=\binom32\,2^2=12,\qquad
N_{\mathrm{faces}}=2^3=8.
\end{equation}%
}
\newcommand{\CurvatureEquilateralExampleDisplay}{%
\begin{equation}
\begin{gathered}
\|\ell e_2-\ell e_1\|=\|\ell e_3-\ell e_1\|
=\|\ell e_3-\ell e_2\|=\sqrt2\,\ell,\\
\frac{\langle\ell e_2-\ell e_1,\ell e_3-\ell e_1\rangle}
{\|\ell e_2-\ell e_1\|\,\|\ell e_3-\ell e_1\|}
=\frac{\ell^2}{2\ell^2}=\frac12=\cos\frac\pi3.
\end{gathered}
\end{equation}%
}
\newcommand{\CurvatureAngleExampleDisplay}{%
\begin{equation}
4\frac\pi3=\frac{4\pi}{3}=240^\circ,\qquad
2\pi-\frac{4\pi}{3}=\frac{2\pi}{3}=120^\circ.
\end{equation}%
}
\newcommand{\CurvatureAngleCountingDisplay}{%
\begin{equation}
\sum_v\sum_{f\ni v}\theta_{vf}
=\sum_f\sum_{v\in f}\theta_{vf}
=\sum_{f=1}^{8}\pi=8\pi.
\end{equation}%
}
\newcommand{\CurvatureScaleExampleDisplay}{%
\begin{equation}
\ell\longmapsto2\ell,\qquad
\sqrt2\,\ell\longmapsto2\sqrt2\,\ell,\qquad
\operatorname{Area}(\tau_s)\longmapsto4\operatorname{Area}(\tau_s),\qquad
\delta_v\longmapsto\delta_v.
\end{equation}%
}

% Worked explanations for Energy.
\newcommand{\EnergyColorBlock}{\Omega_{11}}
\newcommand{\EnergyMixedConnection}{\Omega_{12},\Omega_{21}}
\newcommand{\EnergyBlockDerivative}{D_{\Omega_{11}}}
\newcommand{\EnergyColorConnection}{A}
\newcommand{\EnergyPairWeight}{(g^{-1}\wedge g^{-1})_{AB}}
\newcommand{\EnergyMetricWork}{\Pi_g}
\newcommand{\EnergyMetricExampleDisplay}{%
\begin{equation}
\mathsf g=\begin{pmatrix}-1&0\\0&1\end{pmatrix},\qquad
e=\begin{pmatrix}2&0\\0&3\end{pmatrix},\qquad
g=e^T\mathsf g e=\begin{pmatrix}-4&0\\0&9\end{pmatrix},\qquad
g(v,w)=-4v^0w^0+9v^1w^1.
\end{equation}%
}
\newcommand{\EnergyCurvatureOrderDisplay}{%
\begin{equation}
(\Omega_g\wedge\Omega_g)(\xi,\zeta)
=\Omega_g(\xi)\Omega_g(\zeta)
-\Omega_g(\zeta)\Omega_g(\xi).
\end{equation}%
}
\newcommand{\EnergyColorCorrectionDisplay}{%
\begin{equation}
\begin{aligned}
F_A&=F_{\Omega_{11}}+D_{\Omega_{11}}\Theta+\Theta\wedge\Theta,\\
F_A+\Delta F&=F_{\Omega_{11}}+\Omega_{12}\wedge \Omega_{21}.
\end{aligned}
\end{equation}%
}
\newcommand{\EnergyWeightExampleDisplay}{%
\begin{equation}
g^{\mu\nu}=\operatorname{diag}(-1,1,1,1),\qquad
(g^{-1}\wedge g^{-1})_{(0,1),(0,1)}=-1,\qquad (g^{-1}\wedge g^{-1})_{(1,2),(1,2)}=1.
\end{equation}%
}
\newcommand{\EnergyProjectionExampleDisplay}{%
\begin{equation}
\mathscr F=(N+P)\mathscr F(N+P)
=N\mathscr FN+N\mathscr FP+P\mathscr FN+P\mathscr FP.
\end{equation}%
}
\newcommand{\EnergyRicciMeaningDisplay}{%
\begin{equation}
\operatorname{Ric}_{\mu\nu}=R^\rho{}_{\mu\rho\nu},\qquad
\operatorname{Scal}_g=g^{\mu\nu}\operatorname{Ric}_{\mu\nu}.
\end{equation}%
}
\newcommand{\EnergyInverseMetricDisplay}{%
\begin{equation}
\delta(g^{-1}g)=0
\quad\Longrightarrow\quad
\delta g^{-1}=-g^{-1}(\delta g)g^{-1}.
\end{equation}%
}
\newcommand{\EnergyStressSubstitutionDisplay}{%
\begin{equation}
\begin{aligned}
\operatorname{tr}_g\Pi
&=2\operatorname{Scal}_g-\tfrac13\operatorname{Scal}_g\operatorname{tr}_g g
=\tfrac23\operatorname{Scal}_g,\\
T&=\frac1{2{\kappa}}\bigl(2\operatorname{Ric}-\tfrac13\operatorname{Scal}_g g
-\tfrac23\operatorname{Scal}_g g\bigr)\\
&=\frac1{\kappa}\bigl(\operatorname{Ric}-\tfrac12\operatorname{Scal}_g g\bigr).
\end{aligned}
\end{equation}%
}
\newcommand{\EnergyHomothetyExplanationDisplay}{%
\begin{equation}
\begin{aligned}
g(\nabla_\xi\mathcal E,\zeta)&=g(\nabla_\zeta\mathcal E,\xi),\\
g(\nabla_\xi\mathcal E,\zeta)+g(\xi,\nabla_\zeta\mathcal E)
&=(\mathcal L_{\mathcal E}g)(\xi,\zeta)=4\nu g(\xi,\zeta).
\end{aligned}
\end{equation}%
}
\newcommand{\EnergySkewSquareExampleDisplay}{%
\begin{equation}
\mathsf A=\begin{pmatrix}0&a&b\\-a&0&c\\-b&-c&0\end{pmatrix},\qquad
-\operatorname{tr}(\mathsf A^2)=2(a^2+b^2+c^2)\ge0.
\end{equation}%
}

% Worked explanations for Spectrum.
\newcommand{\SpectrumSurfaceMeasure}{d\Sigma}
\newcommand{\SpectrumNormalizedMap}{U}
\newcommand{\SpectrumConnectionCorrection}{\Delta A}
\newcommand{\SpectrumRegionExampleDisplay}{%
\begin{equation}
\rho|_{\mathcal D}=\rho_0,\qquad
E_{\mathcal D}[n]=\int_{\mathcal D}\rho_0\,d\Sigma
=\rho_0\operatorname{Vol}(\mathcal D).
\end{equation}%
}
\newcommand{\SpectrumIsotropicExampleDisplay}{%
\begin{equation}
r_1=r_2=r_3=r,\qquad \operatorname{Scal}_g=3r,\qquad
\operatorname{Ric}_0=0,\qquad \rho=\frac{3r}{2{\kappa}},\qquad
\mathcal K=\tfrac32r^2=\tfrac23{\kappa}^2\rho^2.
\end{equation}%
}
\newcommand{\SpectrumAnisotropyExampleDisplay}{%
\begin{equation}
\begin{gathered}
(r_1,r_2,r_3)=(r+a,r-a,r),\qquad \operatorname{Scal}_g=3r,\qquad \rho=\frac{3r}{2{\kappa}},\\
|\operatorname{Ric}_0|_h^2=a^2+(-a)^2=2a^2,
\qquad \mathcal K=4a^2+\tfrac32r^2.
\end{gathered}
\end{equation}%
}
\newcommand{\SpectrumSymmetryExampleDisplay}{%
\begin{equation}
\nabla_{(\mu}\xi_{\nu)}=0
\quad\Longrightarrow\quad
\nabla_\mu(T^{\mu\nu}\xi_\nu)=0.
\end{equation}%
}
\newcommand{\SpectrumMassExampleDisplay}{%
\begin{equation}
E=\tfrac54m_*c^2,\qquad
|\mathbf P|=\tfrac34m_*c,\qquad m_*>0,
\qquad M^2=\left(\tfrac{25}{16}-\tfrac9{16}\right)m_*^2=m_*^2.
\end{equation}%
}
\newcommand{\SpectrumNormalizationExampleDisplay}{%
\begin{equation}
\begin{gathered}
h_q=I_3,\qquad \mathsf S=\operatorname{diag}(2,3,4),\qquad
\mathsf G=\mathsf S^\dagger \mathsf S=\operatorname{diag}(4,9,16),\\
\mathsf N=\operatorname{diag}\!\left(\tfrac12,\tfrac13,\tfrac14\right),
\qquad \mathsf N^\dagger\mathsf G\mathsf N=I_3,\qquad
U=\mathsf S\mathsf N=I_3.
\end{gathered}
\end{equation}%
}
\newcommand{\SpectrumCorrectionExampleDisplay}{%
\begin{equation}
D_{A'}(\Delta A)=d(\Delta A)+A'\wedge\Delta A+\Delta A\wedge A'.
\end{equation}%
}


% Worked explanations for Differentiation.
\newcommand{\DifferentiationSecondResponse}{Z}
\newcommand{\DifferentiationBaseDirection}{\eta}
\newcommand{\DifferentiationContactFunction}{\psi(s)=\varphi(\Phi(X(s)))}
\newcommand{\DifferentiationCoefficientDisplay}{%
\begin{equation}
Dw_p(X)[\delta X]=\sum_{j\ne p}\bigl[\langle (\delta X)_p,X_j\rangle_{\mathrm F}+\langle X_p,(\delta X)_j\rangle_{\mathrm F}\bigr].
\end{equation}%
}
\newcommand{\DifferentiationMatrixCurveDisplay}{%
\begin{equation}
\begin{gathered}
X_1(s)=\begin{pmatrix}1&0\\0&s\end{pmatrix},\qquad
X_2(s)=\begin{pmatrix}2&0\\0&1\end{pmatrix},\qquad w_1=w_2=2+s,\\
\Lambda(X(s))_1=\begin{pmatrix}2+s&0\\0&2s+s^2\end{pmatrix},\qquad
\Lambda(X(s))_2=\begin{pmatrix}4+2s&0\\0&2+s\end{pmatrix}.
\end{gathered}
\end{equation}%
}
\newcommand{\DifferentiationMatrixDerivativesDisplay}{%
\begin{equation}
\begin{gathered}
\left.\frac{d}{ds}\Lambda(X(s))_1\right|_0=\begin{pmatrix}1&0\\0&2\end{pmatrix},\qquad
\left.\frac{d}{ds}\Lambda(X(s))_2\right|_0=\begin{pmatrix}2&0\\0&1\end{pmatrix},\\
\left.\frac{d^2}{ds^2}\Lambda(X(s))_1\right|_0=\begin{pmatrix}0&0\\0&2\end{pmatrix},\qquad
\left.\frac{d^2}{ds^2}\Lambda(X(s))_2\right|_0=0.
\end{gathered}
\end{equation}%
}
\newcommand{\DifferentiationChainDefinitionDisplay}{%
\begin{equation}
Y=\Lambda(X),\qquad Z=\Lambda(Y),\qquad {L_X}=D\Lambda(X),\qquad {L_Y}=D\Lambda(Y).
\end{equation}%
}
\newcommand{\DifferentiationNormalizationDisplay}{%
\begin{equation}
v(s)=\begin{pmatrix}1\\s\end{pmatrix},\qquad
{e_r}(s)=\frac{1}{\sqrt{1+s^2}}\begin{pmatrix}1\\s\end{pmatrix},\qquad
{e_r}'(0)=\begin{pmatrix}0\\1\end{pmatrix},\qquad
{e_r}''(0)=\begin{pmatrix}-1\\0\end{pmatrix}.
\end{equation}%
}
\newcommand{\DifferentiationUnitConstraintDisplay}{%
\begin{equation}
\langle {e_r},{e_r}\rangle=1
\quad\Longrightarrow\quad
\langle {e_r},{e_r}'\rangle=0,\qquad
\langle {e_r},{e_r}''\rangle=-\|{e_r}'\|^2.
\end{equation}%
}
\newcommand{\DifferentiationContactTaylorDisplay}{%
\begin{equation}
\psi(0)=0,\qquad
\psi(s)=s \psi'(0)+\tfrac12s^2\psi''(0)+o(s^2),\qquad s\ge0.
\end{equation}%
}
\newcommand{\DifferentiationContactExamplesDisplay}{%
\begin{equation}
\begin{aligned}
\psi(s)=s\quad(s\ge0)&:\qquad \psi'(0)=1,\quad \psi''(0)=0,\\
\psi(s)=s^2\quad(s\in\mathbb R)&:\qquad \psi'(0)=0,\quad \psi''(0)=2.
\end{aligned}
\end{equation}%
}

% Closed coupled source and regional mass displays
\newcommand{\ClosedPatchDirectionDisplay}{%
\begin{equation}
\widehat r_s(\lambda)=\frac{(s_1\lambda_1,s_2\lambda_2,s_3\lambda_3)}{\sqrt{\lambda_1^2+\lambda_2^2+\lambda_3^2}}.
\label{eq:source-patch-direction}
\end{equation}%
}
\newcommand{\ClosedSourcePatchDisplay}{%
\begin{equation}
\KeyMath{
\ExplainTerm{\Psi_s(\lambda;r)}{normalized source patch}
=\ExplainTerm{\left[\frac{R_0}{R(\mathsf D(q_*+r\widehat r_s(\lambda)))}\right]^{1/324}}{fixed-scale normalization}
\ExplainTerm{\mathsf D(q_*+r\widehat r_s(\lambda))}{source on the radial boundary}
}
\label{eq:source-patch}
\end{equation}%
}
\newcommand{\ClosedCurvatureBlockDisplay}{%
\begin{equation}
\Omega_g=\begin{pmatrix}\omega&\theta\\\theta^T&0\end{pmatrix},\qquad
\mathscr F=\begin{pmatrix}F_\omega+\theta\wedge\theta^T&0\\0&0\end{pmatrix},\qquad
F_\omega=d\omega+\omega\wedge\omega.
\label{eq:assembled-curvature-block}
\end{equation}%
}
\newcommand{\ClosedStressMatrixDisplay}{%
\begin{equation}
\KeyMath{
Q=\ExplainTerm{3}{normal-product contribution}+\ExplainTerm{\tfrac12\operatorname{Scal}_\gamma}{spatial scalar contribution},\\
\ExplainTerm{T_{\hat a\hat b}=\frac1{\kappa R^2}\begin{pmatrix}Q&0\\0&-\mathcal D\end{pmatrix}}{energy and stress in the orthonormal rest frame}
}
\label{eq:assembled-stress}
\end{equation}%
}
\newcommand{\ClosedDensityStressDisplay}{%
\begin{equation}
\begin{gathered}
\ExplainTerm{\rho=\frac{Q}{\kappa R^2}}{rest energy density},\qquad
\ExplainTerm{\boldsymbol\pi=0}{rest-frame momentum density},\\
S=\ExplainTerm{-\frac{Q}{3\kappa R^2}I_3}{isotropic stress}
-\ExplainTerm{\frac{\mathcal D_0}{\kappa R^2}}{anisotropic stress},\qquad
\ExplainTerm{\operatorname{tr}_{h_R}S=-\rho}{energy--stress trace relation}
\end{gathered}
\label{eq:assembled-density-stress}
\end{equation}%
}
\newcommand{\ClosedInteriorStressDisplay}{%
\begin{equation}
\KeyMath{
\ExplainTerm{\int_{\Sigma_-}\operatorname{tr}_{h_R}S\,dV_{h_R}}{integrated spatial-stress trace}
=-\ExplainTerm{\int_{\Sigma_-}\rho\,dV_{h_R}}{interior energy}
=-E_-(R,r)>0,\quad0<r<r_0
}
\label{eq:interior-stress-counterpart}
\end{equation}%
}
\newcommand{\ClosedRegionalMassDisplay}{%
\begin{equation}
M^2c^4=(E_-+E_+)^2-c^2|\mathbf P_-+\mathbf P_+|^2.
\label{eq:regional-mass}
\end{equation}%
}
\newcommand{\ClosedSourceMassDisplay}{%
\begin{equation}
\KeyMath{
\ExplainTerm{M}{invariant mass}
=\ExplainTerm{\frac{R}{\kappa c^2}}{scale and physical normalization}
\sum_s\int\ExplainTerm{\chi_s}{chart weight}
\ExplainTerm{\left(3+\tfrac12\operatorname{Scal}_{\gamma,s}\right)}{curvature trace}
\ExplainTerm{\sqrt{\det(\gamma_{ij}^{(s)})}\,d^3q_s}{spatial volume in the chart}
}
\label{eq:source-mass}
\end{equation}%
}
\newcommand{\ClosedInteriorAccountingDisplay}{%
\begin{equation}
E_-=Mc^2-E_+.
\label{eq:interior-energy-accounting}
\end{equation}%
}

% Exact retained source and complete boundary reduction
\newcommand{\SourceBasisDisplay}{%
\begin{equation}
\mathsf D(x,y,z)=\begin{pmatrix}
x&y&-1\\ z&(-y-1)/3&0\\
\zeta&(-y+2)/3&0\\-x-z-\zeta&(-y-1)/3&1
\end{pmatrix},\qquad
\zeta=ax+by-d+(a+b+d)z.
\label{eq:retained-source-basis}
\end{equation}%
}
\newcommand{\SourceReferenceDisplay}{%
\begin{equation}
\mathsf D(q_*)=\begin{pmatrix}-1&-1&-1\\1&0&0\\0&1&0\\0&0&1\end{pmatrix}.
\end{equation}%
}
\newcommand{\SourceDensityCoefficientDisplay}{%
\begin{equation}
\mathfrak f_*=\frac{N_*}{D_*}<0,\quad N_*<0,\quad D_*>0,\qquad
\mathfrak f_*\simeq-1.933636\times10^{39}.
\label{eq:retained-exact-density}
\end{equation}%
}
\newcommand{\SourceInteriorBoundDisplay}{%
\begin{equation}
E_-(R,r)=\frac R{\kappa}\int_{|q-q_*|\leq r}\mathfrak f\,d^3q
\leq\frac{2\pi R}{3{\kappa}}\mathfrak f_*r^3<0,\qquad 0<r<r_0.
\label{eq:source-interior-bound}
\end{equation}%
}
\newcommand{\SourceGramReductionDisplay}{%
\begin{equation}
\mathsf Z=\mathsf Y\mathsf W,\qquad G_Z=\mathsf W G_Y\mathsf W,\qquad \mathsf H=G_Z^2=\mathsf W G_Y\mathsf W^2G_Y\mathsf W.
\end{equation}%
}
\newcommand{\SourceOrderedCurvatureDisplay}{%
\begin{equation}
\begin{gathered}
\mathscr F_{ab}=\ExplainTerm{\mathsf W^{-1}G_Y^{-1}\bigl[(B_a^\perp)^TB_b^\perp-(B_b^\perp)^TB_a^\perp\bigr]\mathsf W}{curvature from the normal response}\\
+\ExplainTerm{\partial_a\Xi_b-\partial_b\Xi_a}{correction derivatives}
+\ExplainTerm{[\beta_a,\Xi_b]-[\beta_b,\Xi_a]}{mixed coupling}
+\ExplainTerm{[\Xi_a,\Xi_b]}{correction commutator}
\end{gathered}
\label{eq:source-ordered-curvature}
\end{equation}%
}
\newcommand{\SourceCurvatureReductionDisplay}{%
\begin{equation}
\mathscr F=
\begin{pmatrix}
F_\omega+\theta\wedge\theta^T&d\theta+\omega\wedge\theta\\
d\theta^T+\theta^T\wedge\omega&\theta^T\wedge\theta
\end{pmatrix}
=\begin{pmatrix}F_\omega+\theta\wedge\theta^T&0\\0&0\end{pmatrix}.
\end{equation}%
}
\newcommand{\SourceRegionalEnergiesDisplay}{%
\begin{equation}
E_-=\frac R{\kappa}\int_{\Sigma_-}Q\,dV_\gamma,\qquad
E_+=\frac R{\kappa}\int_{\Sigma_+}Q\,dV_\gamma.
\label{eq:source-regional-energy}
\end{equation}%
}
\newcommand{\SourceBoundaryReductionDisplay}{%
\begin{equation}
\begin{aligned}
\mathfrak v^k&=\sqrt{\det\gamma}(\gamma^{ij}\Gamma^k{}_{ij}-\gamma^{ki}\Gamma^j{}_{ij}),\\
\mathfrak w&=\sqrt{\det\gamma}\left[3+\tfrac12\gamma^{ij}(\Gamma^k{}_{i\ell}\Gamma^\ell{}_{jk}-\Gamma^k{}_{ij}\Gamma^\ell{}_{k\ell})\right].
\end{aligned}
\label{eq:source-boundary-density}
\end{equation}%
}
\newcommand{\SourceBoundaryIntegralDisplay}{%
\begin{equation}
\int_U \mathfrak f\,d^3q=\int_U\mathfrak w\,d^3q+
\frac12\int_{\partial U}\mathfrak v^k\,dS_k.
\label{eq:source-boundary-integral}
\end{equation}%
}

% Intermediate steps for the appendix derivations.
\newcommand{\SourceDensityScalingDisplay}{%
\begin{equation}
\begin{aligned}
dV_{h_R}&=R^3dV_\gamma=R^3\sqrt{\det\gamma}\,d^3q,
&\rho&=\frac{Q}{{\kappa}R^2},\\
\rho\,dV_{h_R}&=\frac R{\kappa} \mathfrak f\,d^3q.
\end{aligned}
\label{eq:source-density-scaling}
\end{equation}%
}
\newcommand{\SourceCoefficientGeometryDisplay}{%
\begin{equation}
R^2=(\mathbf1^T\mathsf H^{-1}\mathbf1)^{-1},\qquad
\mathbf v=R^2 \mathsf H^{-1}\mathbf1,\qquad
\mathsf g=\mathsf H-2R^2\mathbf1\mathbf1^T.
\label{eq:source-coefficient-geometry}
\end{equation}%
}
\newcommand{\SourceInducedConnectionDisplay}{%
\begin{equation}
\begin{aligned}
P_Y&=\mathsf YG_Y^{-1}\mathsf Y^T,&
B_a^\perp&=(I-P_Y)\partial_a\mathsf Y,\\
\beta_Y&=G_Y^{-1}\mathsf Y^T d\mathsf Y.
\end{aligned}
\label{eq:source-induced-connection}
\end{equation}%
}
\newcommand{\SourceConnectionTransformDisplay}{%
\begin{equation}
\begin{aligned}
\beta&=\mathsf W^{-1}\beta_Y\mathsf W+\mathsf W^{-1}d\mathsf W,\\
F(\beta)&=\mathsf W^{-1}F(\beta_Y)\mathsf W,\\
F(\beta_Y)_{ab}&=G_Y^{-1}
((B_a^\perp)^TB_b^\perp-(B_b^\perp)^TB_a^\perp).
\end{aligned}
\label{eq:source-connection-transform}
\end{equation}%
}
\newcommand{\SourceCurvatureExpansionDisplay}{%
\begin{equation}
\mathscr F=F(\beta)+d\Xi+\beta\wedge\Xi+\Xi\wedge\beta+\Xi\wedge\Xi.
\label{eq:source-curvature-expansion}
\end{equation}%
}
\newcommand{\SourceRicciStressReductionDisplay}{%
\begin{equation}
\begin{aligned}
R_{\hat a\hat b}&=\frac1{R^2}
\begin{pmatrix}0&0\\0&QI_3-\mathcal D\end{pmatrix},
&\operatorname{Scal}_g&=\frac{2Q}{R^2},\\
T_{\hat a\hat b}&=\frac1{\kappa}\left(R_{\hat a\hat b}
-\tfrac12\eta_{\hat a\hat b}\operatorname{Scal}_g\right)
=\frac1{{\kappa}R^2}\begin{pmatrix}Q&0\\0&-\mathcal D\end{pmatrix},
&\eta_{\hat a\hat b}&=\operatorname{diag}(-1,I_3).
\end{aligned}
\label{eq:source-ricci-stress-reduction}
\end{equation}%
}
\newcommand{\SourceObserverWorkDisplay}{%
\begin{equation}
\frac{dE_\alpha}{dR}
=-\int_{\Sigma_\alpha}S^{ij}K_{ij}\,dV_{h_R}
=-\frac1R\int_{\Sigma_\alpha}\operatorname{tr}_{h_R}S\,dV_{h_R}
=\frac{E_\alpha}{R},\qquad \alpha\in\{-,+\}.
\label{eq:source-observer-work}
\end{equation}%
}
\newcommand{\SourceWeightedBoundaryDisplay}{%
\begin{equation}
\int_U\chi \mathfrak f\,d^3q
=\int_U\left(\chi\mathfrak w-\tfrac12\mathfrak v^k\partial_k\chi\right)d^3q
+\frac12\int_{\partial U}\chi\mathfrak v^k\,dS_k.
\label{eq:source-weighted-boundary}
\end{equation}%
}

% Fixed projection and face assignment throughout the manuscript.
\definecolor{carrierred}{HTML}{BE4654}
\definecolor{carriergreen}{HTML}{398268}
\definecolor{carrierblue}{HTML}{4277AD}
\definecolor{diagonalgold}{HTML}{AC8542}
\definecolor{diagonalviolet}{HTML}{8B6A9E}
\tikzset{carrierstar/.style={star,star points=5,star point ratio=2.2,draw=white,line width=.3pt,minimum size=7.8pt,inner sep=0pt},carrieredge/.style={draw=black!55,line width=.65pt},pendingedge/.style={draw=black!18,densely dashed,line width=.6pt},carrierlabel/.style={font=\scriptsize,fill=white,fill opacity=.87,text opacity=1,inner sep=2pt}}
% Arguments: displayed response direction count, edge stage, face stage.
\newcommand{\CarrierScene}[3]{%
\begin{tikzpicture}[x={(1.8cm,-.55cm)},y={(.95cm,.85cm)},z={(0cm,1.65cm)}]
\useasboundingbox (-3.7cm,-3.0cm) rectangle (3.7cm,3.0cm);
\coordinate (rp) at (1,0,0);\coordinate (rm) at (-1,0,0);
\coordinate (gp) at (0,1,0);\coordinate (gm) at (0,-1,0);
\coordinate (bp) at (0,0,1);\coordinate (bm) at (0,0,-1);
\filldraw[fill=carrierred,fill opacity=.035,draw=carrierred!28,line width=.35pt]
 (0,-1.15,-1.15)--(0,1.15,-1.15)--(0,1.15,1.15)--(0,-1.15,1.15)--cycle;
\node[carrierlabel,text=carrierred] at (0,1.12,.72) {$\CarrierDirectionLabel{1}$};
\draw[carrierred!50,line width=.5pt] (-1.1,0,0)--(1.1,0,0);
\ifnum#1=1
 \draw[carrierred!70,-{Stealth[length=4pt]},domain=15:337,samples=70,smooth,variable=\t]
 plot ({0},{.40*cos(\t)},{.40*sin(\t)});
 \node[carrierlabel] at (0,-.55,.32) {$\PhaseRotation$};
\fi
\ifnum#1>1
 \filldraw[fill=carriergreen,fill opacity=.035,draw=carriergreen!28,line width=.35pt]
 (-1.15,0,-1.15)--(1.15,0,-1.15)--(1.15,0,1.15)--(-1.15,0,1.15)--cycle;
 \node[carrierlabel,text=carriergreen] at (-1.12,0,.79) {$\CarrierDirectionLabel{2}$};
 \draw[carriergreen!50,line width=.5pt] (0,-1.1,0)--(0,1.1,0);
\fi
\ifnum#1>2
 \filldraw[fill=carrierblue,fill opacity=.035,draw=carrierblue!28,line width=.35pt]
 (-1.15,-1.15,0)--(1.15,-1.15,0)--(1.15,1.15,0)--(-1.15,1.15,0)--cycle;
 \node[carrierlabel,text=carrierblue] at (1.1,.98,0) {$\CarrierDirectionLabel{3}$};
 \draw[carrierblue!50,line width=.5pt] (0,0,-1.1)--(0,0,1.1);
\fi
\ifnum#3>0
 \fill[carrierred,opacity=.14] (rp)--(gm)--(bm)--cycle;
 \fill[carriergreen,opacity=.14] (rm)--(gp)--(bp)--cycle;
 \fill[carriergreen,opacity=.14] (rm)--(gp)--(bm)--cycle;
 \fill[carrierblue,opacity=.14] (rp)--(gm)--(bp)--cycle;
 \fill[carrierred,opacity=.14] (rp)--(gp)--(bm)--cycle;
 \fill[carrierblue,opacity=.14] (rm)--(gm)--(bp)--cycle;
\fi
\ifnum#3>1
 \fill[diagonalgold,opacity=.17] (rp)--(gp)--(bp)--cycle;
 \fill[diagonalviolet,opacity=.17] (rm)--(gm)--(bm)--cycle;
\fi
\ifnum#2>0
 \foreach \a in {rp,rm}{\foreach \b in {gp,gm}{\draw[carrieredge] (\a)--(\b);}}
\fi
\ifnum#2>1
 \foreach \a in {gp,gm}{\foreach \b in {bp,bm}{\draw[carrieredge] (\a)--(\b);}}
 \ifnum#2=2
  \foreach \a in {rp,rm}{\foreach \b in {bp,bm}{\draw[pendingedge] (\a)--(\b);}}
 \fi
\fi
\ifnum#2>2
 \foreach \a in {rp,rm}{\foreach \b in {bp,bm}{\draw[carrieredge] (\a)--(\b);}}
\fi
\foreach \a in {rp,rm}{\node[carrierstar,fill=carrierred] at (\a) {};}
\node[carrierlabel,anchor=west,xshift=5pt] at (rp) {$\VertexLabel{1}{+}$};
\node[carrierlabel,anchor=east,xshift=-5pt] at (rm) {$\VertexLabel{1}{-}$};
\ifnum#1>1
 \foreach \a in {gp,gm}{\node[carrierstar,fill=carriergreen] at (\a) {};}
 \node[carrierlabel,anchor=south west,xshift=3pt,yshift=3pt] at (gp) {$\VertexLabel{2}{+}$};
 \node[carrierlabel,anchor=north east,xshift=-3pt,yshift=-3pt] at (gm) {$\VertexLabel{2}{-}$};
\fi
\ifnum#1>2
 \foreach \a in {bp,bm}{\node[carrierstar,fill=carrierblue] at (\a) {};}
 \node[carrierlabel,anchor=south,yshift=5pt] at (bp) {$\VertexLabel{3}{+}$};
 \node[carrierlabel,anchor=north,yshift=-5pt] at (bm) {$\VertexLabel{3}{-}$};
\fi
\ifnum#1=2
 \node[carrierlabel] at (0,0,-1.35) {$\FirstPairFigureLabel$};
\fi
\ifnum#1>2
 \node[carrierlabel] at (0,0,-1.4) {$\AdjacentPairsFigureLabel$};
\fi
\ifnum#3>1
 \node[font=\small] at (0,0,-1.75) {$\ColorGroup$};
\fi
\end{tikzpicture}}
\newcommand{\FaceKey}[1]{%
\par\smallskip
\begin{tikzpicture}[x=1cm,y=1cm]
\foreach \x/\y/\col/\lab/\sig in {
0/0/carrierred/{\RedAntiGreen}/{\RedAntiGreenSigns},3.0/0/carriergreen/{\GreenAntiRed}/{\GreenAntiRedSigns},6.0/0/carriergreen/{\GreenAntiBlue}/{\GreenAntiBlueSigns},9.0/0/carrierblue/{\BlueAntiGreen}/{\BlueAntiGreenSigns},
0/-1.05/carrierred/{\RedAntiBlue}/{\RedAntiBlueSigns},3.0/-1.05/carrierblue/{\BlueAntiRed}/{\BlueAntiRedSigns}}
 {\filldraw[fill=\col!12,draw=\col!70] (\x,\y)--(\x+.55,\y+.68)--(\x+1.1,\y)--cycle;
 \node[font=\small] at (\x+.55,\y+.23) {$\lab$};
 \node[font=\scriptsize,anchor=west] at (\x+1.22,\y+.25) {$\sig$};}
\ifnum#1=1
 \foreach \x/\sig in {6.0/{\FirstDiagonalSigns},9.0/{\SecondDiagonalSigns}}
 {\draw[black!35,dashed] (\x,-1.05)--(\x+.55,-.37)--(\x+1.1,-1.05)--cycle;
 \node[font=\scriptsize,text=black!45] at (\x+.55,-.83) {};
 \node[font=\scriptsize,anchor=west] at (\x+1.22,-.8) {$\sig$};}
\else
 \foreach \x/\col/\lab/\sig in {6.0/diagonalgold/{\FirstDiagonalModeLabel}/{\FirstDiagonalSigns},9.0/diagonalviolet/{\SecondDiagonalModeLabel}/{\SecondDiagonalSigns}}
 {\filldraw[fill=\col!14,draw=\col!70] (\x,-1.05)--(\x+.55,-.37)--(\x+1.1,-1.05)--cycle;
 \node[font=\small] at (\x+.55,-.83) {$\lab$};
 \node[font=\scriptsize,anchor=west] at (\x+1.22,-.8) {$\sig$};}
\fi
\end{tikzpicture}}

\title{A Geometric View of Einstein--Yang--Mills Energy and Stress}
\author{Sina Montazeri}
\date{}
\begin{document}
\maketitle
\begin{abstract}
Yang--Mills theory describes interacting gauge fields whose classical Lagrangian contains no gluon mass term. A topological view describes how the radial and transverse directions of a field define oriented phase planes, each carrying a U(1) rotation. Coupling the field and its successive responses links three such planes, generating an SU(3) symmetry whose eight components are represented by the faces of an octahedron.

The resulting closed structure separates its interior from the surrounding spacetime. Its coupled responses determine a spacetime metric and curvature. Variation of the curvature energy yields the stress--energy tensor, including the gravitational response surrounding the structure. Its energy and momentum determine its mass without introducing an elementary gluon mass term into the Yang--Mills Lagrangian.
\end{abstract}
\par\begingroup\centering
\resizebox{.62\textwidth}{!}{\CarrierScene{3}{3}{2}}\par
\smallskip{\small Octahedral representation of the eight $\ColorAlgebra$ components, representing the symmetries of the structure of a gluon in 3 dimensional space.\par}
{\small The arrows $X\to Y\to Z$ denote a field and its successive responses under the coupling $\Lambda$.\par}
\endgroup
\clearpage
\section{Field}
\label{sec:field}
A field assigns values to configurations of a system. The values considered here are matrix components whose magnitudes and relative directions determine their interactions. The physical question is how such interactions can produce the mass of a closed structure when the Yang--Mills Lagrangian contains no elementary gluon mass term.

\subsection{From field interaction to collective mass}
In the usual spacetime description, a gauge connection has matrix-valued components $\FieldYMPotential$, where $\FieldYMCoordinateIndex$ labels a spacetime direction. Using Hermitian matrices, equal to their conjugate transposes, its curvature, or field strength, is
\FieldYMStrengthDisplay
Here $\FieldYMCoupling$ is the coupling strength and
\FieldYMCommutatorDisplay
Directional derivatives measure changes of the connection, while the commutator records the order dependence of its internal transformations. This product term allows the field to interact with itself.

Squaring the field strength retains both propagation and interaction. With $c$ the speed of light and $\hbar$ the reduced Planck constant, choose units where $\FieldYMNaturalUnits$. The Yang--Mills Lagrangian density is
\FieldYMLagrangianDisplay
Repeated indices are contracted with the spacetime metric. Expanding the square gives quadratic, cubic and quartic terms in the connection, with no separate mass term proportional to $\FieldYMMassTerm$~\cite{yangMills1954}. For the gluon field the gauge group is $\ColorGroup$, and the linearized transverse waves are massless. Appendix~\ref{app:coupling} gives the weak-wave calculation.

Absence of an elementary mass term does not preclude collective invariant mass. Restoring $c$, a system with total energy $\FieldYMTotalEnergy$ and total momentum $\FieldYMTotalMomentum$ has invariant mass $\FieldYMCollectiveMass$ satisfying
\FieldYMInvariantMassDisplay
Even massless components can have a nonzero value of this invariant when their momenta cancel. That kinematic observation alone does not establish binding. A glueball is a gauge-invariant bound excitation of the interacting gluon field, with mass determined by its complete rest energy. The weak-wave comparison, a two-excitation example and a lattice result are given in Appendix~\ref{app:coupling}.

The construction follows field comparisons through coupled phase rotations to a closed geometry, then uses its stress and complete energy--momentum to obtain a mass expression. The mass statement retains the condition that the total four-momentum is defined, future-directed and timelike. A common positive threshold for all physical excitations is a distinct mass-gap question~\cite{jaffeWitten}. The first step here is to determine what information about the field components can generate the proposed geometry.

\subsection{Comparing components}
Magnitude alone does not specify an interaction between directions. Two components may have the same length and different alignments. An inner product measures length by self-pairing and alignment by cross-pairing, allowing one direction to be resolved along another.

For example, consider
\FieldExampleVectorsDisplay
Both vectors have unit length,
\FieldExampleLengthsDisplay
but their cross-pairing is
\FieldExampleAlignmentDisplay
Since $\FieldExampleVectorA$ is a unit vector, this value is the component of $\FieldExampleVectorB$ along it. Subtracting that component leaves a perpendicular remainder,
\FieldExampleProjectionDisplay
This is the projection operation used below to obtain radial and transverse field directions.

Matrix components admit the same comparison by multiplying corresponding entries and adding. For two real $\FieldExampleMatrixSize$ matrices,
\FieldExampleMatricesDisplay
the pairing is
\FieldExampleMatrixComparisonDisplay
and self-pairing gives
\FieldExampleMatrixLengthDisplay
Transposition and matrix multiplication collect these products on the diagonal of $A^TB$. Their sum, the trace, is therefore the same inner product. Appendix~\ref{app:coupling} shows the entrywise calculation.

To specify the field components, regard a matrix as a linear map on a finite-dimensional real inner-product space $\FieldVectorSpace$. A map is self-adjoint when its action can be transferred between the two entries of every pairing,
\FieldSelfAdjointDisplay
In an orthonormal basis this condition means that its matrix is symmetric. The numerical example and transpose identity in Appendix~\ref{app:coupling} make this equivalence explicit.

Write $\FieldComponentSpace$ for the real vector space of these self-adjoint components. In orthonormal coordinates its inner product is
\FieldMatrixPairingDisplay
The indices range over the basis directions. Positivity follows from the sum of squared entries, while cross-pairing retains the relative alignment needed for the coupling.

\subsection{The complete field and its variations}
Several labelled components require several positions in the field value. For two components, these positions form
\FieldTwoComponentSpaceDisplay
with a typical value
\FieldTwoComponentValueDisplay
Keeping the positions separate allows the same pairing to compare complete field variations by adding the comparisons at matching labels.

A change of both components is recorded as
\FieldTwoComponentChangeDisplay
and its pairing with $\FieldOtherChangePair$ is
\FieldTwoComponentMetricDisplay
This product metric measures the complete change without losing the identity of either component. A worked matrix example is given in Appendix~\ref{app:coupling}.

For a finite label set $\FieldLabelSet$, let $\mathcal V=\FieldComponentSpace^{\FieldLabelSet}$. On a regular configuration domain $\FieldConfigurationDomain$, the field is
\FieldDefinitionDisplay
and its product metric is
\FieldProductMetricDisplay
Coupling and compatibility determine which field values and changes are admissible; the metric compares their magnitudes and directions. Normalizing a component and projecting another perpendicular to it now applies the elementary comparison to the field itself.

\section{Phase}
\label{sec:phase}
Applying the matrix pairing to a nonzero component gives its radial direction. A second component supplies an independent direction wherever its perpendicular remainder is nonzero. Removing its parallel part gives
\PhaseTransverseDisplay
Normalizing the two directions then yields
\PhaseNormalizationDisplay
Their order specifies an oriented plane. The same operation can first be seen in the vector example from Section~\ref{sec:field}.

\subsection{An oriented plane and its rotations}
Returning to those vectors,
\FieldRadialVectorsDisplay
the radial unit vector is
\FieldRadialUnitDisplay
Subtracting the projection of $\FieldExampleVectorB$ along it gives
\FieldTransverseSubtractionDisplay
The remainder has length $\FieldTransverseLength$, so its normalization is
\FieldTransverseUnitDisplay
Together the two directions satisfy
\FieldPlaneOrthonormalityDisplay
Every vector in their plane therefore has the form
\FieldPlaneVectorDisplay
The first coordinate measures the radial component and the second the transverse component. Ordering these directions makes a positive quarter-turn take $\FieldRadialUnit$ toward $\FieldTransverseUnit$.

Combining the coordinates into
\FieldPlaneComplexCoordinateDisplay
expresses that rotation as multiplication by $\FieldImaginaryUnit$,
\FieldPlaneImaginaryRotationDisplay
which sends the coordinate pair to
\FieldPlaneQuarterTurnDisplay
The quarter-turn defines the complex structure $J$. Applying it twice reverses the vector,
\FieldPlaneComplexStructureDisplay
where $I$ is the identity on this plane.

The matrix field carries the same construction. For example,
\PhaseMatrixExampleDisplay
projection and normalization give
\PhaseMatrixProjectionDisplay
In this basis the quarter-turn acts by
\PhaseComplexStructureDisplay
and rotation through an arbitrary angle is
\PhaseRotationDisplay
A rotation through $\PhaseExampleRotationAngle$ gives
\PhaseMatrixRotationDisplay
The resulting matrix retains unit length while redistributing its radial and transverse components.

Thus a complex coordinate describes two real directions of the same plane. Their inner product becomes
\PhaseHermitianExampleDisplay
A complex line with this inner product is a Hermitian line. Its length-preserving rotations form $\PhaseGroup$. The phase-factor calculation in Appendix~\ref{app:coupling} verifies the arbitrary-angle rotation and composition law.

\begin{figure}[!htbp]
\centering
\resizebox{.62\linewidth}{!}{\CarrierScene{1}{0}{0}}
\caption{Radial and transverse directions span the phase plane. Opposite vertices $\VertexPairOne$ mark the first axis of the representation; the plane carries the internal U(1) rotation.}
\label{fig:phase}
\end{figure}

Figure~\ref{fig:phase} places this first phase direction in the geometric representation. Its display axis records the direction being compared, while the radial and transverse coordinates describe rotation within its internal plane. The next operation is to compare that plane as the field varies.
\FloatBarrier

\subsection{Comparing a moving phase plane}
Let $\FieldPathParameter$ parameterize a path of field configurations, with unit directions $\FieldMovingRadial$ and $\FieldMovingTransverse$. Projecting the radial derivative along the transverse direction measures its internal angular rate,
\FieldConnectionRateDisplay
Differentiating unit length and orthogonality fixes the other in-plane component,
\FieldConnectionProjectedDisplay
Components outside the plane describe its motion in the surrounding matrix space. The projected components specify the rotation that must enter a comparison within the plane.

For several configuration directions, collect these rates in the real connection one-form
\PhasePhaseConnectionDisplay
Here $d$ differentiates along the configuration coordinates. Restricting to a path gives
\PhaseConnectionPathDisplay
and the projected frame derivatives are
\PhaseProjectedDerivativesDisplay
A derivative of vector coordinates must include this frame change as well as their ordinary variation.

Writing a vector in the moving frame gives
\FieldConnectionVectorDisplay
Differentiating both coordinates and basis produces
\FieldConnectionVectorDerivativeDisplay
In complex notation, the same comparison is
\FieldConnectionCovariantDerivativeDisplay
The connection term compensates for basis rotation. Appendix~\ref{app:coupling} verifies this for a fixed vector expressed in a rotating frame.

With local configuration coordinates $x^\mu$, extending the comparison along each coordinate direction gives
\FieldConnectionCoordinateDirectionsDisplay
On a spatial region with coordinates $\FieldSpatialCoordinates$, the same connection has components
\FieldSpatialConnectionDisplay
Here the spatial coordinate $z$ is distinct from the complex coordinate of the internal plane. The spatial labels identify where comparisons are made; the internal plane identifies what rotates in each comparison.

\subsection{Curvature and loop comparison}
Two successive directional comparisons need not agree when their order is exchanged. The covariant derivative and its curvature are
\PhaseUnitaryConnectionDisplay
In two local configuration coordinates, write
\PhaseCurvatureCoordinatesDisplay
The coefficient measures the difference between the two orders of differentiation, per oriented coordinate area. In the spatial description this becomes
\FieldSpatialCurvatureDisplay
The three terms compare small loops in the three coordinate planes for the same phase connection.

Integrating these local comparisons around a loop gives the phase accumulated over its spanning surface. For an oriented surface $\SpanningSurface$ contained in a local patch, Stokes' theorem gives
\PhaseLoopPhaseDisplay
where $\PhaseSurfaceBoundary$ is the oriented boundary. For example,
\PhaseLoopExampleDisplay
represents a quarter-turn. With the convention $\PhaseConnectionConvention$, parallel-transported coordinates satisfy
\PhaseParallelTransportDisplay
Their compensating sign is the same frame-versus-coordinate relation used by the covariant derivative.

A change of phase frame alters the connection coefficients while preserving this curvature. For a smooth phase angle $\PhaseGaugeAngleSymbol$,
\PhasePhaseGaugeDisplay
The added exact differential has vanishing curvature because its mixed derivatives commute. The spatial components transform together,
\FieldSpatialGaugeDisplay
as illustrated by
\PhaseGaugeExampleDisplay
For a single-valued smooth phase angle on the patch, its differential also has zero loop integral. The connection therefore compares the same phase geometry in either frame.

This construction compares one phase plane through changes of the field. When the response depends on the components' mutual pairings, those changes also alter the response. Carrying both together introduces the coupled directions considered next.

\section{Coupling}
\label{sec:coupling}
A phase connection compares internal directions as the field changes. Coupling also changes the response to those directions. On the complete field space $\ComparisonSpace$, the response map is
\CouplingCouplingDisplay
The label $p$ selects the component, while $j$ runs over its partners. Their pairings supply the coefficient multiplying that component, so its response depends on both magnitude and relative alignment.

\subsection{How a change reaches the response}
For a two-component example, take
\CouplingMatrixInputExampleDisplay
Each component has one partner, giving
\CouplingMatrixResponseExampleDisplay
Changing either matrix can therefore affect both responses through their shared pairing.

The first variation must retain that dependence. Writing $L_X=D\Lambda(X)$ gives
\CouplingCouplingResponseDisplay
The first term changes the component at its current coefficient; the second changes the coefficient itself. In the example, varying only the second diagonal entry of the first component yields
\CouplingMatrixVariationExampleDisplay
The second response changes although its own input component is fixed. The derivative $L_X$ records this transmission of a change through the coupling.

\subsection{Two complex response directions}
To compare a field change with its response, place them in $\mathcal V\oplus\mathcal V$ with the product metric of Section~\ref{sec:field}. The graph consists of pairs $(v,L_Xv)$. Its perpendicular directions are obtained from the metric adjoint,
\CouplingAdjointMeaningDisplay
which transfers $L_X$ between the two entries of the pairing. The graph and normal lifts are
\CouplingPairLiftDisplay
Their comparison gives
\CouplingPairOrthogonalityDisplay
so a normal lift is perpendicular to every graph lift. Appendix~\ref{app:coupling} expands the cancellation in coordinates.

Apply both lifts to the radial and transverse directions from Section~\ref{sec:phase}. Let $\PlaneInjection$ insert their linear combinations in the chosen labelled component of $\mathcal V$. Lifting its columns gives
\CouplingPairColumnsDisplay
The graph retains $\TwoColumnFrame$ in its first position and the normal lift retains it in its second. Both planes therefore retain two independent directions, even if $L_X$ has a kernel.

Their Gram matrices satisfy
\CouplingPairGramDisplay
The zero cross-pairing records orthogonality, while the positive Gram matrices give a metric on each plane. Explicitly, for a nonzero plane-coordinate column $\CouplingPlaneCoordinates$,
\CouplingGramLengthsDisplay
Each squared length contains the strictly positive contribution from $\TwoColumnFrame\CouplingPlaneCoordinates$.

Both planes inherit the radial and transverse ordering. Their metrics and orientations therefore supply a quarter-turn on each, giving two orthogonal complex directions. In oriented orthonormal coordinates,
\FieldCouplingComplexPairDisplay
their combined length is
\FieldCouplingComplexNormDisplay
Transformations preserving this complex inner product form $\FieldCouplingUnitaryGroup$.

Preserving the complex volume form imposes the determinant-one condition. For two complex columns,
\CouplingComplexVolumeDisplay
so the unitary maps preserving this form constitute $\PairGroup$. They have matrices
\CouplingPairGroupDisplay
where $\CouplingAlpha,\CouplingBeta$ are complex and the bars denote conjugation. Metric preservation gives orthonormal columns, and determinant one retains their complex volume.

For example,
\CouplingUnitaryMixingExampleDisplay
mixes the two directions equally while retaining unit length and determinant one. Such pair transformations need not commute. Two of them are
\FieldCouplingMixingMatricesDisplay
Applied to
\FieldCouplingInitialCoordinateDisplay
in opposite orders, they give
\FieldCouplingFirstOrderDisplay
and
\FieldCouplingSecondOrderDisplay
The different results exhibit the order dependence that the shared pair actions will carry into the three-direction construction.

\subsection{The shared intermediate response}
Applying the same coupling to $Y$ produces $Z$. The derivative of this next response must pass through the change of $Y$,
\CouplingSuccessiveResponseDisplay
Here $L_Y$ is evaluated at the intermediate field. In the matrix example the next pairing is eight, giving
\CouplingSecondMatrixResponseDisplay
Thus all three levels are determined by the original matrices and the same rule; the two adjacent comparisons share $Y$.

\begin{figure}[!htbp]
\centering
\begin{minipage}[t]{.495\linewidth}\centering
\resizebox{\linewidth}{!}{\CarrierScene{2}{1}{0}}\par\smallskip
\small (a) Two response components and their pair action
\end{minipage}\hfill
\begin{minipage}[t]{.495\linewidth}\centering
\resizebox{\linewidth}{!}{\CarrierScene{3}{2}{0}}\par\smallskip
\small (b) Three coupled response components
\end{minipage}
\caption{Successive responses $X\to Y\to Z$ under $\CouplingMap$ give two adjacent pair actions, shown by solid edges. Dashed edges mark the first--third relation generated from them in Section~\ref{sec:color}.}
\label{fig:couplings}
\end{figure}
\FloatBarrier

The shared intermediate response links the two solid pair relations in Figure~\ref{fig:couplings}. Determining what their combined action generates is the next step.
\FloatBarrier

\section{Algebra}
\label{sec:color}
Retaining the field, its response and the next response gives three linked entries. With $L_X=D\Lambda(X)$ and $L_Y=D\Lambda(Y)$, a variation becomes
\AlgebraThreeLevelVariationDisplay
Both adjacent comparisons use the middle entry. Their pair actions therefore act within a common three-direction construction.

\subsection{Three coupled complex directions}
The graph and normal lifts now lie in $\TripleComparisonSpace$ with its product metric. As for a single pair, the adjoint identity makes each adjacent normal lift perpendicular to the graph. The two normal lifts can overlap through their middle components, so subtracting that overlap is needed to obtain mutually orthogonal planes.

Apply these operations to the same radial and transverse inclusion $\mathsf V$. Write $\mathsf V_0,\mathsf V_1,\mathsf V_2$ for the graph-plane map and the two resulting normal-plane maps. Their pairings satisfy
\AlgebraJointOrthogonalityDisplay
Each map retains two independent columns from $\mathsf V$. Appendix~\ref{app:coupling} gives the projection and the component argument establishing that rank.

The inherited order of the columns orients each plane. Its metric then supplies the quarter-turn constructed in Section~\ref{sec:phase}, so the three planes form a Hermitian three-space. The adjacent actions $\mathrm{SU}(2)_{12}$ and $\mathrm{SU}(2)_{23}$ preserve its metric and complex volume form. Since both act on its middle direction, their generators can also produce a first--third action.

\subsection{Completing the pair symmetry}
To calculate that action, let $E_{ij}$ be the matrix with entry one at $(i,j)$ and zero elsewhere. The trace-free skew-Hermitian generators of a pair are
\AlgebraGeneratorDefinitionDisplay
Here $\mathsf A_{ij}$ and $\mathsf B_{ij}$ mix the two complex coordinates, while $\mathsf H_{ij}$ changes their relative phase. Taking brackets of adjacent generators gives
\AlgebraCrossPairDisplay
Thus the first--third mixing is obtained from the two actions sharing the middle direction. Its diagonal generator satisfies
\AlgebraDiagonalDependenceDisplay

Together these generators supply three complex off-diagonal entries and two independent imaginary diagonal entries,
\AlgebraDimensionDisplay
The bracket calculation supplies the pair relations, and the count identifies their eight independent directions. They span $\ColorAlgebra$, with connected generated group $\ColorGroup$.

\subsection{The connection of the coupled frame}
These directions vary with the source. To compare them at neighboring configurations, choose an oriented orthonormal basis in each plane. The three bases form a real six-frame $\RealFrame$; pairing the radial and transverse entries in each plane expresses the same frame in three complex coordinates.

Let $J$ be the direct sum of their quarter-turns. Differentiating the real frame gives $\Omega_{\mathbb R}=\mathsf E^*d\mathsf E$, where the star denotes the real metric adjoint. Its part commuting with $J$ acts complex-linearly. Expressing that part as a complex $3\times3$ matrix and removing its common phase gives
\AlgebraCompatibleConnectionDisplay
The trace is taken in the three complex coordinates, and $I_3$ is their identity. The first projection preserves the complex structure; the second leaves relative phases.

Orthonormality makes the resulting connection skew-Hermitian, and removal of the common phase makes it trace-free. With $\dagger$ denoting conjugate transpose,
\AlgebraConnectionPropertiesDisplay
It can therefore be written
\AlgebraConnectionMatrixDisplay
The real one-forms $a_i$ describe diagonal phase changes, while $A_{12},A_{23},A_{13}$ mix pairs. All entries depend on the source through its responses, projections and normalization. Appendix~\ref{app:color-correspondence} compares this connection in different normalized color descriptions.

The connection now makes the earlier order dependence local. Its curvature is
\AlgebraGaugeCurvatureDisplay
with first--third entry
\AlgebraCrossCurvatureDisplay
The product $A_{12}\wedge A_{23}$ carries the adjacent interactions into that entry, the curvature counterpart of the bracket that generated the first--third pair.

Multiplying the skew-Hermitian components by $-i$ gives Hermitian coordinates. The two diagonal modes have the normalized basis
\AlgebraDiagonalModesDisplay
The off-diagonal coordinates form three conjugate pairs, completing the eight components. The letters $r,g,b$ name the color basis directions and a bar denotes anticolor.

\begin{figure}[!htbp]
\centering
\begin{minipage}[t]{.495\linewidth}\centering
\resizebox{\linewidth}{!}{\CarrierScene{3}{3}{0}}\par\smallskip
\small (a) All three pair relations
\end{minipage}\hfill
\begin{minipage}[t]{.495\linewidth}\centering
\resizebox{\linewidth}{!}{\CarrierScene{3}{3}{1}}\par\smallskip
\small (b) Six off-diagonal color coordinates
\end{minipage}
\FaceKey{1}
\caption{Adjacent pair actions generate the first--third relation. The six off-diagonal color labels occupy three conjugate pairs of faces; the two remaining faces carry the diagonal modes.}
\label{fig:offdiagonal}
\end{figure}

Figure~\ref{fig:offdiagonal} organizes the components just obtained. Their eight face labels provide the representation from which the closed source surface is assembled.
\FloatBarrier

\section{Closed geometry}
\label{sec:octahedron}
The three complex directions and eight connection components from Section~\ref{sec:color} now give the axes and face labels of an octahedron. Closing this representation supplies the patch arrangement; placing the actual source on those patches then gives the smooth field and its separating boundary.

\subsection{The closed representation}
Let $e_1,e_2,e_3$ be the standard basis of the three-dimensional drawing space and choose
\CurvatureSignedEndpointsDisplay
with drawing scale $\DisplaySize>0$. Taking one signed endpoint from each axis produces
\CurvatureOctahedralFacesDisplay
The eight sign triples give eight faces. Every edge belongs to two faces with opposite induced orientations, so their boundary contributions cancel,
\CurvatureClosedSurfaceDisplay

This closed arrangement also fixes the total angular defect. At a vertex $v$, let $\delta_v$ be a full turn minus the incident face angles. Their sum is
\CurvatureTotalDefectDisplay
Local defects can redistribute under nondegenerate piecewise Euclidean deformations that preserve the connectivity, but their total remains fixed. Appendix~\ref{app:source-reduction} supplies the count. This angular result belongs to the representation; the smooth field's curvature is obtained from its metric after assembly.

\begin{figure}[!htbp]
\centering
\resizebox{.62\linewidth}{!}{\CarrierScene{3}{3}{2}}
\FaceKey{2}
\caption{Closed representation of the eight color components, with six off-diagonal and two diagonal face labels. Shared edges cancel, and the total angular defect is independent of the drawing scale $\DisplaySize$.}
\label{fig:complete}
\end{figure}
\FloatBarrier

\subsection{Assembling the smooth field}
Radial projection turns each face into a patch on the unit sphere. For a sign triple $s=(s_1,s_2,s_3)$, let $\lambda_i\ge0$ locate a point in the face, with $\lambda_1+\lambda_2+\lambda_3=1$. Its unit direction is
\ClosedPatchDirectionDisplay
A vanishing weight places the point on a shared edge, where the adjacent descriptions agree. At a vertex a smooth chart contains the four incident sectors. Together these patches cover the sphere.

To put the constructed field on these patches, use its original four source vectors $b_0,b_1,b_2,b_3$. The columns of $\mathsf D(q)$ record $b_1-b_0,b_2-b_0,b_3-b_0$ in source coordinates $q=(x,y,z)$. Their weighted differences undergo the vector coupling and assemble the self-adjoint matrices $X_p$; the same matrix coupling then gives $Y=\Lambda(X)$ and $Z=\Lambda(Y)$. Thus $\mathsf D$ describes the input to the response chain of Sections~\ref{sec:coupling}--\ref{sec:color}. Appendix~\ref{app:source-coordinates} specifies the weights and assembly.

With reference point $q_*=(-1,-1,1)$, the point $q_*+r\widehat r_s(\lambda)$ selects this source on a sphere of coordinate radius $r>0$. Each patch carries the resulting $X,Y,Z$ and their induced geometry.

These matrices need a common source scale. Let $R(\mathsf D)$ be the positive scale function determined by the coefficient geometry and set $R_0=R(\mathsf D(q_*))$. The normalized patch is
\ClosedSourcePatchDisplay
The exponent follows from $R(a\mathsf D)=a^{324}R(\mathsf D)$ for $a>0$. This factor changes the scale of the matrix while leaving the selected source-coordinate point fixed. Radius $r$ locates the surface; $R_0$ specifies the common scale of its data.

Each patch is a restriction of the same smooth source map. Source values and derivatives therefore agree at shared edges and vertices after the corresponding coordinate changes. Applying the coupling gives a smooth field, metric and connection across the assembled surface.

\subsection{The boundary and its two sides}
On the fixed-scale spatial slice $\Sigma_R$, the image of the coordinate ball $B_r(q_*)$ defines $\Sigma_-$. Its boundary is $S_r=\partial\Sigma_-$, and the complementary region is $\Sigma_+=\Sigma_R\setminus\overline{\Sigma_-}$. The signs label sides of the surface. The radius remains general within the regular neighborhood of the source.

Allowing the source coordinates and scale to vary continues this geometry across $S_r$. Neighborhoods on its two sides, called collars, have matching induced metrics and second fundamental forms when compared with a common normal. Thus the separating surface divides a matched geometry with source data available on both sides.

\FloatBarrier

\section{Energy and stress}
\label{sec:energy}
Energy and momentum are read from the response of this geometry to source variations. The calculation begins with the matrix pairings that determine its metric.

\subsection{Metric and complete curvature}
Collect the terminal matrices as columns $\mathsf Z=(\operatorname{vec}Z_p)_p$. Their Frobenius pairings give $G_Z=\mathsf Z^T\mathsf Z$, and $\mathsf H=G_Z^2$ supplies a positive metric on their coefficients. On the regular source family, the scale is $R^2=(\mathbf1^T\mathsf H^{-1}\mathbf1)^{-1}$ and the normalized coefficient vector is $\mathbf v=R^2\mathsf H^{-1}\mathbf1$, where $\mathbf1$ is the column of ones.

Reversing the metric sign along this scale direction gives $\mathsf g=\mathsf H-2R^2\mathbf1\mathbf1^T$. The coframe $e$ is the covariant change of $\mathbf v$ along the source history, computed with the $\mathsf H$-compatible connection. It maps a tangent vector to its coefficient change, so pairing two such images gives the spacetime metric. With $\Omega_g$ its torsion-free metric connection,
\NativeGeometryDisplay
Appendix~\ref{app:source-metric} derives these coefficients and the coframe from the same response matrices. Changing to spacetime coordinates introduces the frame derivative $de$ in $\Gamma(g)$. Differentiation then gives the curvature in the two descriptions,
\NativeCurvatureDisplay
Thus $\mathscr F$ and $\mathcal R$ represent the same curvature before and after the frame change.

The color connection $A$ from Section~\ref{sec:color} enters $\Omega_g$ with the mixed components and connection corrections. Let $\mathscr F_{\mu\nu}$ denote its curvature two-form components for $\mu,\nu=0,1,2,3$, and $R_{abcd}$ the Riemann tensor in an orthonormal frame. Contracting all contributions gives
\NativeRiemannEnergyDisplay
The metric raises indices with their Lorentzian signs. Appendix~\ref{app:source-reduction} expands the weights and blocks, and Appendix~\ref{app:color-correspondence} gives the color-frame transformation. Varying this complete invariant now determines how the geometry responds to a source change.

\subsection{From curvature work to physical stress}
One contraction of Riemann curvature gives Ricci curvature; its trace gives the scalar curvature,
\EnergyRicciMeaningDisplay
The trace-free remainder is the Weyl tensor $\mathcal C$. In four dimensions the invariant separates into
\NativeWeylEnergyDisplay
Here $\langle\,\cdot\,,\,\cdot\,\rangle_g$ contracts all tensor indices using $g^{-1}$. This tensor contraction uses the generated metric, while the matrix pairing in Section~\ref{sec:field} compared the original field components. The scalar term corrects for the trace already included in the Ricci term.

A source variation changes both curvature and the metric used to contract it. Write $\Pi$ for the coefficient of $\delta\operatorname{Ric}$ and $\Pi_g$ for that of $\delta g$ in the Ricci and scalar contributions. Including the Weyl contribution gives
\NativeFullWorkDisplay
The first term varies the Weyl tensor and its metric pairing. The other terms record the Ricci and metric responses, with frame and coupling coefficients changing along the same source family. Their coefficient calculation is given in Appendix~\ref{app:source-reduction}.

The physical work measurement follows the convergence and separation of neighboring source-response curves. After retaining their acceleration and frame transport, these focusing measurements reconstruct Ricci curvature on the regular source family. Its conjugate work coefficient is therefore the same $\Pi=2\operatorname{Ric}-\tfrac13\operatorname{Scal}_g\,g$ obtained in the variation.

The calibration fixes the relative normalization of trace-free and scalar focusing deformations, while $\kappa>0$ converts to physical units. It preserves their work pairing, as shown in Appendix~\ref{app:curvature-work}. In this convention the coefficient paired with the calibrated deformation is
\NativeStressExtractionDisplay
Substitution of $\Pi$, with $\operatorname{tr}_g g=4$, yields
\NativeEinsteinDisplay
Thus the measured focusing stress agrees with the Einstein tensor divided by $\kappa$. The full source variation still contains the Weyl and metric work displayed above.

The resulting stress also supplies the curvature source current. Let $D_\mu$ be the derivative induced by $\Omega_g$ on coefficient-frame endomorphisms, $\nabla_\mu$ the Levi-Civita derivative of $g$, and $j_\nu$ the source current obtained from the antisymmetrized Ricci derivative. The field and stress balances are
\NativeSourcedEquationsDisplay
Appendix~\ref{app:source-reduction} gives the current and frame conversion. These equations apply to the matched geometry and its continuation beyond the boundary.

Substituting the stress relation back into the invariant gives
\NativeEnergyStressDisplay
This expresses the same complete contraction through stress and Weyl curvature. To obtain a measured density, the tensor $T$ must next be resolved along the source's observer and spatial directions.

\subsection{Energy and stress of the closed source}
For the homogeneous family, $g=-dR^2+R^2\gamma$, where $\gamma$ is the spatial metric at unit scale. Choose the future normal $n=\partial_R$; the slice metric is $h_R=R^2\gamma$. Let $\theta=(\theta^i)$ be an orthonormal coframe for $\gamma$, with dual vectors $E_i$ and connection $\omega$. Placing the spatial directions before the normal in the connection gives
\ClosedCurvatureBlockDisplay
Torsion freedom cancels the mixed curvature blocks. The product $\theta\wedge\theta^T$ remains in the spatial curvature and contributes on each of its three oriented two-planes.

Represent that curvature on $(E_2\wedge E_3,E_3\wedge E_1,E_1\wedge E_2)$ by the symmetric matrix $\mathcal D$, with trace $Q$. For the physical reading, use the orthonormal rest frame $(n,R^{-1}E_1,R^{-1}E_2,R^{-1}E_3)$; hatted indices denote this ordering.

In this frame, $\rho=T(n,n)$ is the energy density. The spatial projections $\pi_i=-T(n,R^{-1}E_i)$ give momentum, and the restriction of $T$ to the slice gives the spatial stress $S$. Ricci contraction and calibrated trace reversal yield
\ClosedStressMatrixDisplay
Each spatial two-plane contributes to the three in $Q$. Trace reversal supplies the energy entry even though the normal Ricci component vanishes. The intermediate Ricci matrix is displayed in Appendix~\ref{app:source-reduction}.

Separating $\mathcal D=(Q/3)I_3+\mathcal D_0$, with $\operatorname{tr}\mathcal D_0=0$, resolves the matrix into
\ClosedDensityStressDisplay
The displayed spatial matrix contains the components of $S$ in the $h_R$-orthonormal frame. Its trace is therefore $-\rho$ for isotropic and anisotropic stress alike. The same reduction gives zero rest-frame momentum. Appendix~\ref{app:source-reduction} relates these quantities to the spatial curvature contraction.

For the reference source, $\mathfrak f=Q\sqrt{\det\gamma}$ is negative at $q_*=(-1,-1,1)$. Since $\rho\,dV_{h_R}=(R/\kappa)\mathfrak f\,d^3q$, continuity gives $r_0>0$ for which the regional integral $E_-(R,r)=\int_{\Sigma_-}\rho\,dV_{h_R}$ is negative when $0<r<r_0$. Integrating the stress trace on the same region yields
\ClosedInteriorStressDisplay
Thus the negative regional energy has an exact positive integrated spatial-stress counterpart.
\FloatBarrier

\section{Collective mass}
\label{sec:collective-mass}
Integrating those densities gives the regional energy and momentum. Their combined balance retains the surrounding response in the invariant mass.

\subsection{Measured energy and momentum}
Let $u$ be a future unit timelike observer and $n$ the future unit normal to a spatial slice $\Sigma$. Contracting $T$ with these directions gives
\NativeObservedEnergyDisplay
The slice metric supplies $d\Sigma$, while $T$ retains the normalization $\kappa$. For the rest observer $u=n$ this contraction is
\NativeObservedDensityDisplay
Homogeneity gives $\operatorname{Ric}(n,n)=0$, recovering $\rho=Q/(\kappa R^2)$ from the preceding stress matrix. Integrating over $\Sigma_-$ and $\Sigma_+$ gives $E_-$ and $E_+$.

Momentum uses the same source-frame comparison. The spatial projections $\pi_i=-T(n,R^{-1}E_i)$ give regional components $P_{\pm,i}=c^{-1}\int_{\Sigma_\pm}\pi_i\,dV_{h_R}$. The factor $1/c$ converts the stress--energy projection to momentum units. Their vanishing in the source rest frame gives $\mathbf P_\pm=0$.

The continued tensor determines $E_+$ and describes the surrounding gravitational response. Equation~\eqref{eq:interior-stress-counterpart} instead integrates the spatial-stress trace over $\Sigma_-$, giving $-E_-$. These two contractions retain the roles established in Section~\ref{sec:energy}.

\subsection{Combining the two regions}
Matching at $S_r$ gives the same stress on both sides, while the regions have opposite outward normals. Their shared-interface fluxes therefore cancel when the regional balances are added. The remaining balance includes the exterior contribution and any actual outer flux.

To track this balance for a changing observation, let $\xi$ be the spacetime observation vector. Stress conservation and the product rule give
\NativeChargeBalanceDisplay
Parentheses symmetrize indices with weight $1/2$. Thus $\nabla_{(\mu}\xi_{\nu)}=(\nabla_\mu\xi_\nu+\nabla_\nu\xi_\mu)/2$ measures the metric deformation along $\xi$. Its contraction with $T$ is the associated work. Integrating retains this work, outer-boundary flux and boundary transport, as detailed in Appendix~\ref{app:source-reduction}. A metric symmetry removes the deformation term; a constant charge also requires zero net outer flux.

Moving $S_r$ at fixed field and fixed $R$ transfers the same contribution between $\Sigma_-$ and $\Sigma_+$, leaving their sum unchanged. The radius $r$ specifies this partition. Changing $R$ changes the source geometry and its energy, with the corresponding work already included in the balance.

\subsection{Mass of the closed structure}
Combining the regional measurements gives $E=E_-+E_+$ and $\mathbf P=\mathbf P_-+\mathbf P_+$. When this total four-momentum is defined, future-directed and timelike, its invariant mass satisfies
\EnergyTotalMassRelationDisplay
Substitution gives
\ClosedRegionalMassDisplay
The cross terms belong to the invariant of the complete system. In the source rest frame $\mathbf P=0$, so $E=Mc^2$ and
\ClosedInteriorAccountingDisplay

The source geometry enters this mass explicitly through its density and volume. Substituting $\rho=Q/(\kappa R^2)$ and $dV_{h_R}=R^3dV_\gamma$ combines the scale factors into $R/\kappa$. The remaining integrand is the curvature trace $Q=3+\tfrac12\operatorname{Scal}_\gamma$ times spatial volume.

To express the integral on source charts $q_s$ covering both regions, take nonnegative weights $\chi_s$ with $\sum_s\chi_s=1$. Each chart has $dV_{\gamma,s}=\sqrt{\det(\gamma_{ij}^{(s)})}\,d^3q_s$, and the weights count overlaps once. Appendix~\ref{app:source-reduction} retains their derivative and outer-boundary contributions in the equivalent bulk--boundary expression.

Under the stated four-momentum conditions, the coupled curvature and spatial volume therefore give the following mass. The source scale $R$ remains explicit, while the fixed-field sum is independent of the partition radius $r$,
\ClosedSourceMassDisplay

\clearpage
\appendix
\section{Comparisons, variations and supporting calculations}
\label{app:coupling}
The calculations here support the field comparison, phase connection and coupled response in Sections~\ref{sec:field}--\ref{sec:color}. First, explicit matrix examples verify the comparisons used in the construction. The later derivatives follow a source variation through its responses and compatibility conditions. A curve parameter $s$ labels these variations and is independent of physical time.

\subsection{Matrix comparisons}
For the two matrices used to explain the field pairing, transposition exchanges rows and columns,
\FieldExampleTransposeDisplay
Multiplication then collects matching column products on the diagonal,
\FieldExampleDiagonalDisplay
and their sum gives
\FieldExampleTraceDisplay
This verifies the trace formula used in Section~\ref{sec:field}.

A self-adjoint map transfers its action between the entries of the inner product. For example,
\FieldSelfAdjointExampleDisplay
gives
\FieldSelfAdjointActionsDisplay
and hence
\FieldSelfAdjointComparisonsDisplay
The equality holds for every vector pair because the matrix is symmetric,
\FieldSelfAdjointIdentityDisplay
This identifies symmetric matrices as the self-adjoint component values used by the field.

The product metric applies the same comparison at every component label. Consider two complete field changes with components
\FieldProductFirstExampleDisplay
and
\FieldProductSecondExampleDisplay
Their component pairings are
\FieldProductComparisonsDisplay
so the complete comparison is
\FieldProductTotalDisplay
Self-pairing similarly gives
\FieldProductLengthDisplay
Each labelled change contributes to the length of the complete variation.

\subsection{Phase rotation and moving frames}
In the oriented plane of Section~\ref{sec:phase}, rotation through $\FieldPlaneAngle$ multiplies its complex coordinate by
\FieldPlanePhaseFactorDisplay
The real coordinates consequently transform as
\FieldPlaneRotationCoordinatesDisplay
Since
\FieldPlanePhaseLengthDisplay
the rotation preserves length. These factors constitute
\FieldPlanePhaseGroupDisplay
and compose by addition of angles,
\FieldPlanePhaseCompositionDisplay
This is the phase action carried by the radial and transverse plane.

To verify the connection's role, keep a plane fixed and rotate its frame,
\FieldConnectionRotatingFrameDisplay
Its angular rate is
\FieldConnectionRotatingRateDisplay
Now take the fixed vector
\FieldConnectionFixedVectorDisplay
Its coordinates in that rotating frame are
\FieldConnectionFixedCoordinatesDisplay
or, in complex form,
\FieldConnectionFixedComplexDisplay
Their ordinary derivative is nonzero, but the frame term cancels it,
\FieldConnectionCancellationDisplay
The covariant derivative therefore records the fixed geometric vector correctly.

An additional frame angle $\FieldGaugeAngle$ changes the basis by
\FieldConnectionFrameChangeDisplay
The coordinate of the same vector changes oppositely,
\FieldConnectionCoordinateChangeDisplay
while the connection becomes
\FieldConnectionPotentialChangeDisplay
Together these transformations give
\FieldConnectionCovarianceDisplay
so either frame describes the same derivative. Choosing $\FieldGaugeUnrotation$ removes the earlier frame rotation and gives constant coordinates for the fixed vector.

\subsection{Weak waves and collective mass}
For the conventional gluon connection, linearization about the trivial connection gives the transverse wave equation
\FieldYMWaveDisplay
Here $z$ denotes a spatial coordinate. With angular frequency $\FieldYMFrequency$ and wave vector $\FieldYMWaveVector$, it implies
\FieldYMDispersionDisplay
Using $\FieldYMQuantumEnergy$ and $\FieldYMQuantumMomentum$, where $\hbar$ is the reduced Planck constant, gives
\FieldYMMasslessEnergyDisplay
Comparison with
\FieldYMRelativisticEnergyDisplay
identifies $m=0$ for these weak-field excitations~\cite{jaffeWitten}.

Their zero individual masses do not require zero mass for a collection. If two excitations have energy $\FieldYMExampleEnergy$ and opposite momenta,
\FieldYMOppositeMomentaDisplay
the totals are
\FieldYMCombinedEnergyDisplay
and their invariant mass is
\FieldYMPairMassDisplay
This illustrates the distinction between constituent mass and collective mass; binding remains a separate dynamical property.

In pure gauge theory, a glueball is a bound excitation of the interacting gluon field. Lattice calculations estimate such collective energies. Morningstar and Peardon obtained approximately
\FieldYMGlueballMassDisplay
for the lightest scalar state, using their stated conversion to physical units~\cite{morningstarPeardon1999}. The label $\FieldYMScalarLabel$ denotes zero spin and positive parity and charge-conjugation parity. This result is a comparison, not an input to the constructed source's mass.

A common mass gap would require a positive threshold for all physical excitations. The opposite-momentum example alone supplies no such threshold because its input energy can be arbitrarily small. The main construction instead follows the specific closed geometry to its complete energy--momentum.

\subsection{Derivatives of the coupling}
The response used in Section~\ref{sec:coupling} is
\FieldCouplingResponseDisplay
with component rule
\FieldCouplingWeightsDisplay
Both factors depend on the source. Writing
\FieldCouplingDerivativeDisplay
gives
\FieldCouplingVariationDisplay
The two terms vary the component and its pairings with the other components.

Writing the same response directly as a polynomial in its input gives
\CouplingMapPolynomialCouplingDisplay
Its differential is the preceding variation with $L_X=D\Lambda(X)$ written out,
\CouplingMapCouplingDifferentialDisplay
The derivative of its scalar coefficient is
\DifferentiationCoefficientDisplay
Each pairing varies through both entries, so an unchanged input component can still acquire a changed response.

Differentiating once more in independent directions $\delta_1X,\delta_2X$ gives
\CouplingMapCouplingSecondDifferentialDisplay
The first two terms differentiate the two factors of the first variation; the last differentiates the pairings again. Exchange of the variation directions leaves the result unchanged.

For the matrix example in Section~\ref{sec:coupling}, let the first component vary in its second diagonal entry,
\DifferentiationMatrixCurveDisplay
The coefficient $w_1=w_2=2+s$ varies with that entry, giving at $s=0$
\DifferentiationMatrixDerivativesDisplay
The quadratic term comes from their product. This verifies why coefficient variation must remain in the response derivative.

\subsection{Successive responses}
To carry the variation into the next level, write
\DifferentiationChainDefinitionDisplay
The derivatives $L_X$ and $L_Y$ are evaluated at different inputs. Since the second input is itself a response, its change must also be differentiated.

For a configuration direction $\DifferentiationBaseDirection$, put $\delta X=DX[\eta]$. The chain rule gives
\CouplingMapFullChainRuleDisplay
The second response receives the composed change $L_YL_X\delta X$, while the last expression differentiates the evaluation point of $L_Y$.

If the varied vector also depends on the configuration, the full derivative is
\CouplingMapBothDifferentialVariationsDisplay
All three contributions enter the changing frame, along with derivatives of its projections and normalization. Treating either response derivative as fixed would omit part of that change.

\subsection{Normalization and compatibility}
Normalization itself contributes nonlinear derivatives. For example,
\DifferentiationNormalizationDisplay
The unnormalized vector has zero second derivative, but its unit direction does not. Differentiating unit length in general gives
\DifferentiationUnitConstraintDisplay
The first identity makes velocity tangent to the unit sphere; the second fixes the acceleration's normal component.

Compatibility must be differentiated along the same source curve. Let $\CombinedFieldMap$ collect the responses, metric, connection and derivatives in a smooth condition $\CompatibilityEquality$. Take a twice differentiable curve $\CompatibleFieldCurve$ satisfying it, and set $\FieldVelocityAcceleration$. Dots refer to $s$, with derivatives of $\Phi$ evaluated at $X(0)$ and those of $f$ at $\Phi(X(0))$.

One and two differentiations give
\CouplingMapPropagatedCompatibilityDisplay
The second-order condition contains source acceleration, nonlinear response and curvature of the constraint. Retaining all three terms ensures that the varied induced quantities follow the source compatibility.

For a smooth inequality $\CompatibilityInequality$ that is an equality initially, a compatible curve with $s\ge0$ obeys
\CouplingMapContactDerivativeDisplay
Writing $\DifferentiationContactFunction$ gives the Taylor expansion
\DifferentiationContactTaylorDisplay
A negative first derivative violates the inequality for sufficiently small positive $s$. If the first derivative vanishes, the second must be nonnegative. A curve allowed for both signs requires zero first derivative.

The distinction is illustrated by
\DifferentiationContactExamplesDisplay
The linear example is admissible on one side and the quadratic example on both. These differential restrictions are necessary; the curve must still satisfy the original compatibility relations.

\subsection{Graph and normal planes}
The first response places a variation in the graph
\FieldCouplingGraphDisplay
The adjoint relation
\FieldCouplingAdjointDisplay
then gives a normal lift
\FieldCouplingNormalDisplay
whose complete pairing with the graph is
\FieldCouplingOrthogonalityDisplay
The cancellation establishes the perpendicular directions used in Section~\ref{sec:coupling}.

Inserting the radial and transverse frame in the selected component gives
\FieldCouplingGraphPlaneDisplay
and
\FieldCouplingNormalPlaneDisplay
The graph retains the two independent directions in its first position; the normal lift retains them in its second. Their ranks therefore remain two.

For an explicit check, take
\FieldCouplingExampleDerivativeDisplay
The paired directions are
\FieldCouplingExampleLiftsDisplay
and, for example,
\FieldCouplingExampleCancellationDisplay
All cross-plane pairings vanish similarly. Each vector has squared length five, so division by $\FieldCouplingExampleNormalization$ gives an orthonormal frame within either plane.

For the three-level joint variation of Section~\ref{sec:color}, the lifts become
\AlgebraTripleLiftDisplay
The continuing matrix example gives
\AlgebraResponseExampleDisplay
for the final graph component. Adjacent adjoint cancellations show
\AlgebraNormalCancellationDisplay
Both normal planes are thus perpendicular to the graph, though they may overlap each other through the middle position.

Subtracting that overlap gives
\AlgebraOrthogonalColumnsDisplay
Here $\mathsf W_2$ is the second normal map before projection. For a plane-coordinate column $\xi\in\mathbb R^2$, the coefficients $\alpha\in\mathbb R^2$ solve
\AlgebraProjectionDisplay
The inverse exists because $\mathsf V_1$ has independent columns. Subtraction leaves the third block of $\mathsf V_2$ equal to $\mathsf V$, since $\mathsf V_1$ has zero third block. All three maps therefore retain two independent columns and the mutual orthogonality used in the main text.

\subsection{Pair-generator calculations}
The matrix units act by
\AlgebraMatrixUnitExampleDisplay
so their product connects the first and third coordinates. For the real pair generator,
\AlgebraGeneratorActionDisplay
multiplication in both orders gives
\AlgebraBracketProductsDisplay
This verifies the first--third generator in Section~\ref{sec:color}. The imaginary generator follows from the same matrix-unit product.

For the diagonal components, the trace pairing gives
\AlgebraDiagonalNormExampleDisplay
and the first--third difference is
\AlgebraDiagonalExpansionDisplay
These relations fix their normalization and dependence.

Finally, connection curvature retains the order of its matrix-valued factors. For local tangent directions $\eta,\zeta$,
\AlgebraWedgeExampleDisplay
The nonzero example illustrates the adjacent interaction in the first--third curvature entry.

\clearpage
\section{Source geometry and energy balance}
\label{app:source-reduction}
This appendix follows the source used in Section~\ref{sec:octahedron} into the curvature and stress of Section~\ref{sec:energy}, then into the regional balance used for mass. The coordinate radius $r$ specifies a dividing surface; the source scale $R$ controls the homogeneous geometry.

\subsection{Angular curvature of the octahedral representation}
The signed vertices in Section~\ref{sec:octahedron} give the convex hull
\CurvatureOctahedralHullDisplay
A point within a face has the form
\CurvatureFaceAverageDisplay
For example,
\CurvatureFacePairExampleDisplay
share the edge joining their first two vertices. With consistent outward face orientations, their induced edge directions satisfy
\CurvatureEdgeCancellationExampleDisplay
Pairing every shared edge this way closes the representation.

Counting the vertices, edges and faces gives
\CurvatureCountingExampleDisplay
Equal axis lengths make each triangular face equilateral,
\CurvatureEquilateralExampleDisplay
For a vertex $v$, let $\theta_{vf}$ be its corner angle in an incident face $f$. The angular defect is
\CurvatureAngularDefectDisplay
Four faces meet there, so
\CurvatureRegularDefectDisplay
or, in degrees,
\CurvatureAngleExampleDisplay

To sum the defects, count all corner angles once by vertex and once by face,
\CurvatureAngleCountingDisplay
Each triangular face contributes $\pi$, giving
\CurvatureInvariantDefectDisplay
This is the total in equation~\eqref{eq:octahedral-total-defect}. Rescaling the drawing preserves it,
\CurvatureScaleExampleDisplay
The calculation concerns the representation's angles. The radial patch map in Section~\ref{sec:octahedron} supplies the smooth surface on which the coupled source is evaluated.

\subsection{Source coordinates}
\label{app:source-coordinates}
That patch map uses an affine source plane with basis matrix $B\in\mathbb Q^{12\times4}$. Write $q=(x,y,z)$ and $\bar q=(x,y,-1,z)^T$. Stacking the rows of $\mathsf D(q)$ into a column gives $B\bar q$, equivalently
\SourceBasisDisplay
The constants $a,b,d$ are entries $B_{7,1},B_{7,2},B_{7,3}$, with indices starting at one. Their exact rational values are supplied in the \href{retained-source-exact-data.json}{supplementary data}; the formula determines the remaining basis entries.

The columns are the differences $b_i-b_0$ for $i=1,2,3$ between four original source vectors in $\mathbb R^4$, with the common translation fixed by $b_0=0$. They generate the matrix field used in the main construction. For each anchor $p\in\{0,1,2,3\}$, use the fixed labels $(\phi_0,\phi_1,\phi_2,\phi_3)=(0,1,3,7)$ and the original positive vector pairing $\beta$ to form
\[
v_{pi}=(\phi_i-\phi_p)(b_i-b_p),\qquad
(\lambda_p x)_i=\left(\sum_{j\ne p,i}\beta(x_i,x_j)\right)x_i,\qquad i\ne p.
\]
Applying this vector coupling twice gives $z_{p\bullet}=\lambda_p(\lambda_p(v_{p\bullet}))$. The self-adjoint matrix
\[
X_p=\sum_{i\ne p}z_{pi}z_{pi}^{\flat_\beta}
\]
is assembled from these responses; $z^{\flat_\beta}(w)=\beta(z,w)$. The matrix-level responses are then exactly $Y=\Lambda(X)$ and $Z=\Lambda(Y)$ from Section~\ref{sec:coupling}. Vector and matrix operations are thus successive levels of the same source construction.

Fixing $\bar q_3=-1$ gives three coordinates transverse to scaling. At $q_*=(-1,-1,1)$,
\SourceReferenceDisplay
Equation~\eqref{eq:source-patch} normalizes this family to a fixed positive scale $R$. The ball $|q-q_*|<r$ then selects the interior region on that slice, distinguishing the role of radius from source scale.

\subsection{Matrix response, metric and induced connection}
\label{app:source-metric}
To calculate the geometry from the coupling, collect $Y=\Lambda(X)$ into columns $\mathsf Y=(\operatorname{vec}Y_p)_p$ and $Z=\Lambda(Y)$ into $\mathsf Z=(\operatorname{vec}Z_p)_p$. Their Gram matrices are $G_Y=\mathsf Y^T\mathsf Y$ and $G_Z=\mathsf Z^T\mathsf Z$. The coupling multiplies column $p$ by $w_p(Y)=\sum_{j\ne p}(G_Y)_{pj}$. With $\mathsf W=\operatorname{diag}(w_p(Y))$,
\SourceGramReductionDisplay
A Gram matrix introduces a weight on each side, and squaring it produces $\mathsf H$. These formulas apply on the regular patches where the columns are independent and $\mathsf W$ is invertible.

The resulting coefficient geometry is
\SourceCoefficientGeometryDisplay
Here $\mathbf1$ is the column of ones. The definition of $R^2$ gives $\mathbf1^T\mathbf v=1$, while $\mathsf H\mathbf v=R^2\mathbf1$. Thus $\mathbf v$ is $\mathsf H$-orthogonal to coefficient directions whose entries sum to zero. The signed metric keeps their positive pairing and gives $\mathsf g(\mathbf v,\mathbf v)=-R^2$ on the scale line.

The coframe is also determined by these matrices. On history coordinates $x^\mu$, let $\mathsf C=\mathsf Z$ and $\mathsf K=G_Z$, so $\mathsf H=\mathsf K^2$. The induced connection and its positive metric correction are
\[
\mathsf b_\mu=\mathsf K^{-1}\mathsf C^T\partial_\mu\mathsf C,\qquad
\mathsf A^+_\mu=\mathsf b_\mu+
\tfrac12\mathsf K^{-1}(\mathsf b_\mu+\mathsf b_\mu^T)\mathsf K.
\]
They satisfy $(\mathsf A^+_\mu)^T\mathsf H+\mathsf H\mathsf A^+_\mu=\partial_\mu\mathsf H$. Covariantly differentiating the normalized coefficient vector gives
\[
e_\mu=\partial_\mu\mathbf v+\mathsf A^+_\mu\mathbf v,\qquad
e=[e_0\ e_1\ e_2\ e_3],\qquad
g_{\mu\nu}=e_\mu^T\mathsf g e_\nu.
\]
The coframe therefore records how source changes alter the normalized coefficient vector, including the moving coefficient metric. It has rank four on the retained regular history. Pairing its columns gives the spacetime metric used in Section~\ref{sec:energy}; homogeneity reduces it to $g=-dR^2+R^2\gamma$.

Differentiating the column space requires its tangential projection and normal remainder,
\SourceInducedConnectionDisplay
In these source-coordinate expressions, $\partial_a=\partial/\partial q^a$ with $a=1,2,3$. Thus $B_a^\perp$ measures the part of the changing columns perpendicular to their current span, while $\beta_Y$ compares directions within that span.

Passing from $\mathsf Y$ to the weighted response $\mathsf Z$ transforms the connection and its curvature,
\SourceConnectionTransformDisplay
The term containing $d\mathsf W$ retains the varying coupling weights. Here $F(\beta)=d\beta+\beta\wedge\beta$, with the same convention for $\beta_Y$. The source-coordinate indices $a,b$ in curvature components are distinct from the fixed constants in the affine source matrix.

\subsection{Complete curvature and calibrated work}
\label{app:curvature-work}
The metric and connection above determine the full curvature before any homogeneous block reduction. Its invariant uses all coordinate components. For increasing pairs $\NativePairIndices$, its weights are
\NativeEnergyDefinitionDisplay
The pair labels collect antisymmetric coordinate indices, each ranging over four spacetime directions.

A two-direction example verifies how the metric transforms with its frame,
\EnergyMetricExampleDisplay
Each scale occurs twice because the metric pairs two vectors. The connection product likewise retains both orders,
\EnergyCurvatureOrderDisplay
so its contraction must preserve the complete matrix action.

Let $N$ be the $\mathsf g$-orthogonal projector onto the timelike normal line and $P=I-N$ its spatial complement. Inserting both projectors gives
\EnergyProjectionExampleDisplay
Grouping the blocks defines
\NativePacketsDisplay
Their cross-pairings vanish under the trace, giving
\NativeEnergyAllocationDisplay
Both the normal--mixed and mixed--tangent descriptions contain $\mathcal K_C$, so one copy is subtracted in their sum. This accounts for all blocks once.

A frame separating color and complementary directions has
\SpectrumFullCurvatureDisplay
Here the two-digit subscripts label blocks. Writing $D_{\Omega_{11}}$ for the derivative induced by the first block, the difference between its connection and $A$ gives
\SpectrumMixedCorrectionDisplay
and
\EnergyColorCorrectionDisplay
The mixed products thus remain before restricting to the homogeneous solution.

To vary the invariant, first differentiate its inverse metric,
\EnergyInverseMetricDisplay
Applying this rule to the Ricci and scalar terms in equation~\eqref{eq:native-weyl-energy} gives
\NativeWorkCoefficientsDisplay
The coefficient $\Pi$ pairs with $\delta\operatorname{Ric}$ and $\Pi_g$ with $\delta g$. The full variation also includes the Weyl term. The physical focusing measurement is obtained from the response of the same source, with acceleration and frame transport retained. For its response velocity $v$, write $\mathsf B^a{}_b=\nabla_bv^a$, $\vartheta=\operatorname{tr}\mathsf B$ and $a_v=\nabla_vv$. The measured directional contraction is
\[
r_{\mathrm{foc}}(v)=\operatorname{div}_g a_v-v(\vartheta)-\operatorname{tr}(\mathsf B^2)
=\operatorname{Ric}(v,v).
\]
The available source probes reconstruct the corresponding symmetric tensor. This is the focusing tensor whose conjugate work coefficient is $\Pi$; the identification is made after the source response and its transport are included.

For a focusing variation $B$, the adopted deformation calibration is $4B-\tfrac43g\operatorname{tr}_gB$. Its pairing with the physical stress obeys
\[
\frac{\kappa}{2}
\left\langle T,\,4B-\frac43g\operatorname{tr}_gB\right\rangle_g
=\langle\Pi,B\rangle_g,\qquad
T=\frac1{2\kappa}\left(\Pi-g\operatorname{tr}_g\Pi\right).
\]
Substitution verifies this work identity using $\operatorname{tr}_g g=4$. The calibration fixes the relative trace and trace-free response; $\kappa$ retains the physical unit conversion. It is an operational normalization supplied with the source-work correspondence, not a new consequence of choosing a tensor norm.

Its conserved trace combination and unit trace-free gain give the displayed stress. In the class of metric-natural linear trace adjustments, the contracted Bianchi identity fixes the trace ratio when scalar curvature varies, while the gain fixes the remaining normalized factor. Keeping $\kappa$ explicit, the substitution is
\EnergyStressSubstitutionDisplay

The same stress relation can be substituted back into the invariant. Trace reversal gives
\NativeTraceReversalNormDisplay
which, together with $\operatorname{tr}_gT=-\operatorname{Scal}_g/\kappa$, proves equation~\eqref{eq:native-energy-stress}.

Antisymmetrizing the covariant derivative of Ricci and transforming to the coefficient frame gives
\NativeCurrentSourceDisplay
The resulting $j_\nu$ is the source current in Section~\ref{sec:energy}, with the same connection and metric.

\subsection{Homogeneous curvature and rest-frame stress}
\label{app:source-stress}
To evaluate the calibrated stress for the retained source, pass from the induced connection $\beta$ to the torsion-free metric connection $\Omega_g$. Here $\beta$ denotes the induced connection, distinct from the vector pairing in the source assembly. Their difference $\Xi=\Omega_g-\beta$ gives
\SourceCurvatureExpansionDisplay
Using the normal derivatives of the column space makes this curvature explicit,
\SourceOrderedCurvatureDisplay
The first two products compare the two orders of normal response. The remaining terms differentiate the correction and retain its mixed and quadratic products. All contribute when the full curvature is contracted.

For the homogeneous family, the adapted connection lists spatial directions before the timelike normal. Torsion freedom, $d\theta+\omega\wedge\theta=0$, gives
\SourceCurvatureReductionDisplay
Both mixed blocks vanish; the bottom scalar block vanishes because each $\theta^i\wedge\theta^i=0$. The product $\theta\wedge\theta^T$ remains in the spatial block.

Represent this spatial curvature on $(E_2\wedge E_3,E_3\wedge E_1,E_1\wedge E_2)$ by $\mathcal D$, where $E_i$ are dual to $\theta^i$. Its trace is $Q=3+\operatorname{Scal}_\gamma/2$. In the physical orthonormal frame $(n,R^{-1}E_1,R^{-1}E_2,R^{-1}E_3)$, the Ricci and stress contractions become
\SourceRicciStressReductionDisplay
Although the normal Ricci entry is zero, scalar trace reversal gives the energy entry $Q/(\kappa R^2)$. In the spatial block the two terms proportional to $QI_3$ cancel, leaving $-\mathcal D/(\kappa R^2)$. Taking that block's trace gives $-\rho$, the relation integrated in Section~\ref{sec:energy}.

\subsection{Reference density and the interior sign}
\label{app:source-density}
The scale enters the slice metric as $h_R=R^2\gamma$. Combining its volume element with the measured density gives
\SourceDensityScalingDisplay
The product contains one factor of $R$, since volume contributes $R^3$ and curvature contributes $R^{-2}$. Defining $\mathfrak f=Q\sqrt{\det\gamma}$ therefore writes the regional energy directly in source coordinates.

At the reference point, the exact value $\mathfrak f_*=\mathfrak f(q_*)$ is rational,
\SourceDensityCoefficientDisplay
The integers $N_*<0$ and $D_*>0$ are given in the supplementary data; the displayed decimal is rounded. Continuity on a regular neighborhood gives $r_0>0$ with $\mathfrak f(q)\le\mathfrak f_*/2$ for $|q-q_*|\le r<r_0$. Integrating over the ball of coordinate volume $4\pi r^3/3$ yields
\SourceInteriorBoundDisplay
The metric volume density is already in $\mathfrak f$. This is the negative regional energy used in Section~\ref{sec:energy}; the exterior remains part of the complete accounting.

\subsection{Curvature contraction on the homogeneous geometry}
Homogeneity explains why curvature has no scale-direction component on this family. Let $\NativeEuler$ be the Euler field generating scaling with degree $\NativeDegree$. Its metric relations are
\NativeHomothetyDisplay
The Lie derivative fixes the metric's scale change, while the exact one-form identifies the Euler field's metric pairing. Its negative norm makes the direction timelike.

Because the pairing is exact, its exterior derivative vanishes. For the torsion-free connection, the lowered covariant derivative is therefore symmetric. Combining this symmetry with the scale relation gives
\EnergyHomothetyExplanationDisplay
In the next formula, $X,Y$ denote local tangent vector fields, not the coupled response levels. The resulting derivative implies
\NativeEulerCurvatureDisplay
The parenthesis is the torsion, so curvature annihilates the Euler direction. Riemann symmetries extend this vanishing to every slot containing it.

Normalize the Euler field to $\NativeEulerNormal$ and complete it to an orthonormal frame. All surviving slots are spatial, where the metric is positive. Hence
\NativeSpatialEnergyDisplay
For a real skew matrix,
\EnergySkewSquareExampleDisplay
the negative trace of its square is twice the sum of its independent squared entries. Each surviving curvature matrix has this form, giving the displayed nonnegative contraction.

To separate its density and anisotropy, write $\NativeSpatialMetric$ and $\NativeSpatialTracefreeRicci$. Here $h=h_R$ is the spatial metric on the chosen slice, distinct from the Hermitian-line pairing used earlier. In a spatial Ricci eigenframe, let $\NativeRicciEigenvalues$ be its eigenvalues and $\NativeSectionalValues$ the sectional curvatures. Their contractions give
\NativeRicciSectionalDisplay
Solving for the sectional values yields
\NativeSectionalSolutionDisplay
and substitution in the full contraction gives
\NativeSpatialRicciEnergyDisplay
Resolving Ricci into trace and trace-free parts, with $\rho=\operatorname{Scal}_g/(2\kappa)$, therefore gives
\NativeDensityIdentityDisplay
The density occurs squared, so this norm identity alone does not select its sign.

For example, equal spatial eigenvalues give
\SpectrumIsotropicExampleDisplay
Here $r$ is their common value, not the dividing radius. Changing two eigenvalues by opposite amounts gives
\SpectrumAnisotropyExampleDisplay
The trace and density stay fixed while the anisotropic contribution changes. This shows what the curvature norm records beyond measured density.

For a constant density, integration gives
\SpectrumRegionExampleDisplay
Here $\mathcal D$ is a region, distinct from the curvature matrix used above. An energy--momentum example gives
\SpectrumMassExampleDisplay
illustrating why invariant mass retains both entries. Neither example supplies a numerical value for the constructed source.

\subsection{Regional balance}
The constructed surface bounds $\Sigma_-$ on the fixed-scale slice; its complement is $\Sigma_+$. Their restrictions of the metric, connection, curvature and stress agree at $S_r$, and their outward normals are opposite. The regional energies are
\SourceRegionalEnergiesDisplay
Adding their balances cancels shared-interface fluxes. Moving $S_r$ within a fixed collecting region transfers a common energy between the two integrals without changing their sum.

Changing $R$ is a different operation. For a region transported at fixed normalized source coordinates, the slice has second fundamental form $K_{ij}=h_{R,ij}/R$, and its deformation work is
\SourceObserverWorkDisplay
The final equality uses $\operatorname{tr}_{h_R}S=-\rho$. A moving outer boundary adds its transport contribution. This is the same observer-deformation term retained in equation~\eqref{eq:native-charge-balance}.

\newpage
\subsection{Coordinate bulk and boundary terms}
To expose the boundary contribution of the density integral, let $\Gamma^k{}_{ij}$ be the Levi-Civita connection of $\gamma$. The scalar-curvature density can be organized as
\SourceBoundaryReductionDisplay
Thus $\mathfrak f=\mathfrak w+\tfrac12\partial_k\mathfrak v^k$ separates a bulk density from a coordinate divergence. For a regular coordinate region $U$,
\SourceBoundaryIntegralDisplay
The oriented element $dS_k$ is a coordinate surface element; $\mathfrak v^k$ already contains $\sqrt{\det\gamma}$.

Weighting a chart by $\chi$ also differentiates that weight,
\SourceWeightedBoundaryDisplay
These overlap terms accompany the weights in the final mass expression. Summing the regional contributions cancels the shared interface and retains the actual outer boundary. This is a decomposition of the already-defined regional energy, with each contribution counted once and the same total-four-momentum qualifications retained.

\clearpage
\section{Color-frame correspondence}
\label{app:color-correspondence}
The color connection used in Sections~\ref{sec:color} and \ref{sec:energy} can be compared with a connection on the configuration quotient. The comparison must retain the metric, normalization and changing frame on the stated regular domains.

Start with the real six-frame from Section~\ref{sec:color}. Differentiating its orthonormality gives
\AlgebraFrameSkewDisplay
so its connection is skew-symmetric. This is the real-frame identity underlying the compatible color projection.

To compare normalized color descriptions, let $\SpectrumGraphMap$ map the differential graph to the configuration quotient. Write $\SpectrumQuotientMetric$ for its positive Hermitian metric matrix and $\SpectrumGramMatrix$ for the graph Gram matrix. These represent the metrics on the corresponding spaces in the stated common dimension.

Let $\SpectrumNormalizer$ convert the graph coordinates to an orthonormal basis, and write $\SpectrumNormalizedMap$ for the resulting map between normalized coordinates. Their metric relations are
\SpectrumActualStateMapDisplay
The first equality says that pulling back the quotient pairing gives the graph pairing. The matrix $\SpectrumNormalizer$ normalizes the graph coordinates. Multiplication by the positive square root of the quotient metric then gives $\SpectrumNormalizedMap$, whose adjoint product is the identity. It therefore preserves the inner product of those normalized color coordinates.

For diagonal metrics the operations are explicit,
\SpectrumNormalizationExampleDisplay
The normalizer removes the graph's three scale factors before the map is applied. This example illustrates the metric equality without changing its domain or assumptions.

As the source varies, its normalized frame varies too. Writing $\SpectrumQuotientConnection$ for the quotient color connection, with $(q)$ denoting this description, gives
\SpectrumActualConnectionDisplay
The derivative of the normalized map accounts for that frame change. Including it makes curvature transform by conjugation.

The connection entering the full field may also differ from this transported connection. Keeping their difference gives
\SpectrumActualCurvatureDisplay
where $\SpectrumConnectionDifference$. Its covariant derivative contains the interaction with the transported connection,
\SpectrumCorrectionExampleDisplay
and the remaining quadratic term contains the correction's self-product. These are the terms retained when the full curvature enters the energy contraction in Section~\ref{sec:energy}.


\clearpage
\bibliographystyle{unsrt}
\bibliography{references}
\end{document}
